%%
%% memoria.tex
%%
\documentclass[a4paper,12pt]{scrbook}
\pagestyle{headings}

% Paquete para dibujos
\usepackage{tikz}

% Referencias externas
\usepackage{xr-hyper}

% Paquete para generar enlaces clickables
\usepackage[pdftex, breaklinks=false, colorlinks=true, linkcolor=black, anchorcolor=black, urlcolor=blue, citecolor=red]{hyperref}

% Encoding y tratamiento de fuentes
\usepackage[T1]{fontenc}
\usepackage[spanish]{babel}
\usepackage[utf8]{inputenc}

% Paquete con el entorno tabularx, que define un tipo de celda X que hace word-wrapping
\usepackage{tabularx}

% Paquete con el entorno longtable y longtabu para tablas de varias páginas
\usepackage{longtable}

% Provee la flag H para forzar los floats en el sitio donde los pongas
\usepackage{float}

% Provee el comando \captionof para dar etiquetas a elementos no flotados
\usepackage{caption}

% Paquete que provee el comando \FloatBarrier para poner una barrera para floats
\usepackage{placeins}

% Paquete para controlar el interlineado a nivel de entorno
\usepackage{setspace}

% Paquete para la fuente charter
\usepackage{charter}

% Paquete para la fuente helvética
\usepackage[scaled=0.92]{helvet}

% Paquete para GNUPlot
\usepackage{gnuplottex}

% Guiones de hyphenado
\usepackage{hyphenat}

% Extra ToC listings
\usepackage{tocbibind}

% Gráficos
\usepackage{graphicx}

% Para poner figuras una al lado de otra
\usepackage[lofdepth,lotdepth]{subfig}

% Paquete para figuras giradas
\usepackage{rotating}

%%%%%%%%%%%%%%%%%%%%%%%%%%%%%%%%%%%%%%%%%%%%%%%%%%%%%%%%%%%%%%%%%%%%%%%%%%%%%%%%%%%%%%%%%%%%%%%%%%%%%%%%%%%%%%
% Paquetes en local (directorio 'paquetes')

% Fuente para códigos
\usepackage{paquetes/inconsolata}

% Automatiza el resaltado de sintaxis usando pygments
\usepackage{paquetes/minted}

% Soporte de acrónimos
\usepackage[printonlyused]{paquetes/acronym-custom}

% Cambia el nombre del apéndice de bibliografía
\addto\captionsspanish{
\renewcommand\bibname{Bibliografía y referencias}
}

%%%%%%%%%%%%%%%%%%%%%%%%%%%%%%%%%%%%%%%%%%%%%%%%%%%%%%%%%%%%%%%%%%%%%%%%%%%%%%%%%%%%%%%%%%%%%%%%%%%%%%%%%%%%%%

\begin{document}

% Órdenes LaTeX variadas
% http://www.tex.ac.uk/cgi-bin/texfaq2html?label=chngmargonfly

\newenvironment{changemargin}[2]{%
  \begin{list}{}{%
    \setlength{\topsep}{0pt}%
    \setlength{\leftmargin}{#1}%
    \setlength{\rightmargin}{#2}%
    \setlength{\listparindent}{\parindent}%
    \setlength{\itemindent}{\parindent}%
    \setlength{\parsep}{\parskip}%
  }%
  \item[]}{\end{list}}

\newenvironment{nota}{
  \begin{changemargin}{2em}{2em}
    \textbf{\textsc{Nota: }}
}{
  \end{changemargin}
}

\newcommand{\nombreProyecto}{SiteUp: plataforma para la vigilancia de la disponibilidad de servicios de Internet}
\newcommand{\jugador}{\textit{Jugador}}
\newcommand{\sistema}{\textit{Sistema}}


% %% PRODUCTOS
% \newcommand{\nombrepostprocesador}{ACL2:\colonhyp{}Procesador}
% \newcommand{\nombrevisor}{XMLEye}
% \newcommand{\nombreyaxml}{YAXML:\colonhyp{}Reverse}
% \newcommand{\postprocesador}{\texttt{\nombrepostprocesador}\xspace}
% \newcommand{\visor}{\nombrevisor\xspace}
% \newcommand{\yaxml}{\texttt{\nombreyaxml}\xspace}

% \newcommand{\biblioteca}[1]{\index{#1}\texttt{#1}}

% % YAML/YAXML
% \newcommand{\etiqueta}[1]{{\bfseries \ttfamily #1}}

% % ACL2

% \newcommand{\orden}[1]{\texttt{#1}}   % Nombre de orden ACL2
% \newcommand{\fichero}[1]{\texttt{#1}} % Nombre de fichero
% \newcommand{\evento}[1]{\texttt{#1}}  % Nombre de un evento
% \newcommand{\libro}[1]{\textsc{#1}}   % Libro ACL2

% % PERL

% \newcommand{\modulo}[1]{\index{módulo Perl!#1}\index{#1|see{módulo Perl!#1}}\texttt{#1}}
% \newcommand{\funcion}[1]{\textit{#1}}

% % JAVA

% \newcommand{\clase}[1]{\textit{#1}}   % Clase Java
% \newcommand{\metodo}[1]{\texttt{#1}}  % Método (también Perl)
% \newcommand{\paquete}[1]{\texttt{#1}}

% %% MANUALES

% \newcommand{\accesoteclado}[1]{\textsc{#1}} % Acceso de teclado

% %% OTROS

% \newcommand{\patron}[1]{\emph{#1}}

% \newcolumntype{,}{>{$}r<{$}}
% \newcommand{\Index}[1]{#1\emph{\index{#1}}}

% \newcounter{pasoacept}

% \newenvironment{pruebaaceptacion}{
%   \setcounter{pasoacept}{0}
%   \begin{center}
%   \begin{tabular}{| >{\stepcounter{pasoacept}\arabic{pasoacept}. }p{.4\linewidth} | p{.5\linewidth}|}
%     \hline
%     \multicolumn{1}{| c |}{\textbf{Paso seguido}} & \multicolumn{1}{c|}{\textbf{Resultado esperado}} \\
%     \hline
%     \hline
% }{
%   \hline
%   \end{tabular}
%   \end{center}
% }

% \renewcommand{\lstlistlistingname}{Listados}
% \renewcommand{\lstlistingname}{Listado}

% \newcommand{\CPP}
% {\mbox{C\hspace{-.1em}\raise.2ex\hbox{+\hspace{-.1em}+}}\xspace}

% \setlength{\extrarowheight}{4pt}

% \lstset{
%   extendedchars,
%   flexiblecolumns,
%   stringstyle=\ttfamily,
%   showstringspaces=false,
%   frame=tb
% }

% %% PARTE DE DOCBOOK

% \newcommand{\application}[1]{\index{#1}\emph{#1}}
% \newcommand{\cmdsynopsis}[1]{\nohyphens{\texttt{#1}}}
% \newcommand*{\command}[1]{\nohyphens{\textbf{\texttt{#1}}}}
% \newcommand{\constant}[1]{\texttt{#1}}
% \newcommand{\computeroutput}[1]{#1}
% \newcommand*{\email}[1]{\nohyphens{\texttt{#1}}}
% \newcommand*{\envar}[1]{\nohyphens{\texttt{#1}}}
% \newcommand*{\filename}[1]{\texttt{#1}}
% \newcommand*{\guibutton}[2][]{\emph{#2}}
% \newcommand*{\guilabel}[2][]{\emph{#2}}
% \newcommand*{\guimenuitem}[2][]{\emph{#2}}
% \newcommand*{\guimenu}[2][]{\emph{#2}}
% \newcommand*{\keycombo}[1]{\textsc{#1}}
% \newcommand*{\keysym}[1]{\textsc{#1}}
% \newcommand*{\option}[1]{\nohyphens{\texttt{#1}}}
% \newcommand{\prompt}[1]{#1}
% \newcommand{\varname}[1]{\nohyphens{\texttt{#1}}}

% \newcommand{\note}[1]{\vskip 1em
%     \fbox{\parbox{\textwidth}{\textsc{Nota}
%     \vskip 1em #1}} \vskip 1em}


%%% Local Variables:
%%% mode: latex
%%% TeX-master: "memoria"
%%% End:


% Portada
\begin{titlepage}
  \centering
  \includegraphics[width=.3\textwidth]{0_misc/logo_uca}

  \bigskip
  \bigskip
  \bigskip

  \begin{changemargin}{2em}{2em}
    \centering

    {\Huge \textsc{\nohyphens{Escuela Superior de Ingeniería}}}

    \bigskip
    \bigskip
    \bigskip

    {\huge \nohyphens{Ingeniería Informática}}

    \bigskip
    \bigskip
    \bigskip
    \bigskip
    \bigskip
    \bigskip

    \begin{doublespace}
      {\LARGE \nohyphens{\nombreProyecto}}
    \end{doublespace}


    \bigskip
    \bigskip
    \bigskip
    \bigskip

    \bigskip
    \bigskip
    \bigskip
    \bigskip
    \bigskip
    \bigskip
    \bigskip

  \end{changemargin}

  {\Large José Tomás Tocino García \\}

  \bigskip

  {\large Cádiz, \today}

\end{titlepage}


%%% Local Variables: 
%%% mode: latex
%%% TeX-master: "../memoria"
%%% End: 

\cleardoublepage

% Segunda portada
{
  \thispagestyle{empty}
  \centering
  \includegraphics[width=.2\textwidth]{0_misc/logo_uca}

  \bigskip
  \bigskip
  \bigskip
  
  \begin{changemargin}{2em}{2em}

    \begin{center}
      {\Huge \textsc{\nohyphens{Escuela Superior de Ingeniería}}}
      
      \bigskip
      \bigskip
      
      {\huge \nohyphens{Ingeniería Informática}}
      
      \bigskip
      \bigskip
      \bigskip
      \bigskip
      
      \begin{doublespace}
        {\LARGE \nohyphens{\nombreProyecto}}
      \end{doublespace}


      \bigskip
      \bigskip
      \bigskip
      \bigskip
      
    \end{center}
  \end{changemargin}
  \begin{changemargin}{2em}{1em}
  \begin{flushleft}
    \Large

    \textsc{Departamento}: \nohyphens{Ingeniería Informática.} \\
    \textsc{Director del proyecto}: \nohyphens{Iván Ruiz Rube} \\
    \textsc{Autor del Proyecto}: \nohyphens{José Tomás Tocino García}. \\
  \end{flushleft}

  \end{changemargin}  

  \bigskip
  \bigskip
  \bigskip
  
  \begin{flushright}
    \large
    Cádiz, \today

    \bigskip
    \bigskip
    \bigskip
    \bigskip    
    \bigskip
    \bigskip

    Fdo.: José Tomás Tocino García
    
  \end{flushright}

}

%%% Local Variables: 
%%% mode: latex
%%% TeX-master: "../memoria"
%%% End: 



\cleardoublepage

% Texto de licencia (pequeño)
\bigskip
\bigskip

Este documento se halla bajo la licencia \ac{FDL}. Según estipula la
licencia, se muestra aquí el aviso de copyright. Se ha usado la
versión inglesa de la licencia, al ser la única reconocida
oficialmente por la \ac{FSF}.

\begin{quote}
  Copyright \copyright~\the\year~José Tomás Tocino García.

  Permission is granted to copy, distribute and/or modify this document
  under the terms of the GNU Free Documentation License, Version 1.2
  or any later version published by the Free Software Foundation;
  with no Invariant Sections, no Front-Cover Texts, and no Back-Cover Texts.
  A copy of the license is included in the section entitled "GNU
  Free Documentation License".
\end{quote}

%%% Local Variables: 
%%% mode: latex
%%% TeX-master: "../memoria"
%%% End: 

\cleardoublepage

% Agradecimientos
\section*{Agradecimientos}
Quiero agradecer y dedicar la presente memoria y, por extensión, el proyecto
SiteUp al completo:
% \begin{itemize}
% \item A Julian Raschke por crear y mantener Gosu.
% \item A Manuel Palomo Duarte y Rafael Rodríguez Galván, de la OSLUCA, por confiar en mí
%   como becario de la oficina, apoyarme y darme facilidades para desarrollar el
%   proyecto durante la beca.
% \item A Antonio García Domínguez, como co-tutor del proyecto, por sus acertados
%   consejos, su visión para descubrir cosas mejorables y su implicación con el
%   software y el conocimiento libre, gracias a la cual esta memoria fue posible.
% \item A mis compañeros más cercanos de la carrera, que me han apoyado y de los
%   que me he nutrido para desarrollar este proyecto.
% \item A mi familia y mi pareja, que finalmente no tendrán que recibir como
%   excusa el desarrollo del proyecto a la hora de dedicarles el tiempo que se
%   merecen.
% \item A Robert Penner por su magnífico conjunto de ecuaciones para animaciones,
%   sin las que oFlute tendría una interfaz mucho más pobre.
% \item A los que apuestan por el software libre y el conocimiento abierto. Sus
%   herramientas, bibliotecas y documentos han hecho posible que oFlute llegase a
%   buen puerto.
%\end{itemize}


\cleardoublepage

% Índices varios
\tableofcontents
\listoffigures
%\listoftables

\setlength{\parskip}{1.2ex plus 0.4ex minus 0.1ex}

\chapter{Introducción}
\label{chap:introduccion}
\section{Contexto y motivación}
Las tecnologías de la información en general e Internet en particular son ya
parte integral de la sociedad. Casi todos los ámbitos de la vida diaria de las
personas, desde las interacciones sociales hasta la búsqueda de empleo, cuentan
ya con su reflejo en las tecnologías de la información, a menudo mediante el uso
de servicios web a través de Internet. 

No solo los aspectos tradicionales de la sociedad tienen su presencia en las
redes, también han surgido nuevos modelos empresariales propios de Internet que
han crecido de manera importante y se han situado a niveles comparables a los de
las empresas tradicionales. Empresas puramente digitales como Facebook o Twitter
ya cotizan en bolsa y realizan operaciones bursátiles del orden de miles de
millones de dólares.

Se pone pues de manifiesto la importancia de la fiablidad de los servicios e
infraestructuras de los que dependen estos nuevos modelos de negocio. El
\textit{uptime} -- un término inglés que describe el porcentaje de tiempo que un
servicio se mantiene disponible -- debe ser siempre cercano al 100\%, dado que
en caso contrario los potenciales usuarios del servicio se encontrarán con que
no pueden acceder a él, dando lugar incluso a pérdidas económicas. Es el caso de
Amazon, que llegó a perder 4.8 millones de dólares al sufrir un fallo que dejó
inaccesible su web durante 40 minutos~\cite{amazon}.

De todo lo expuesto se extrae la necesidad de contar con sistemas para
monitorizar la disponibilidad de estos servicios y, en caso necesario, actuar de
manera que puedan subsanarse las causas de los problemas.

\section{Objetivos}
\label{sec:objetivos}

A la hora de definir los objetivos de un sistema, podemos agruparlos
en dos tipos diferentes: \textbf{funcionales} y
\textbf{transversales}. Los primeros se refieren a \textit{qué} debe
hacer la aplicación que vamos a desarrollar, e inciden
directamente en la experiencia del usuario y de potenciales
desarrolladores.

Por otro lado, los objetivos transversales son aquellos invisibles al
usuario final, pero que de forma inherente actúan sobre el resultado
final de la aplicación y sobre la experiencia de desarrollo de la misma.

\subsection{Funcionales}
\begin{itemize}
\item Crear un conjunto de herramientas para la monitorización y el chequeo de
  diversos aspectos del estado de un servicio de Internet.
\item Crear una aplicación online, de acceso público, que permita la creación y
  gestión de chequeos de manera sencilla, basada internamente en las
  herramientas mencionadas en el punto anterior.
\item Habilitar esta aplicación de un sistema de notificaciones mediante correo
  electrónico que alerte a los usuarios de posibles cambios en la disponibilidad
  de los servicios monitorizados.
\item Crear una aplicación móvil para el sistema operativo Android para la
  recepción instantánea de avisos provenientes de la aplicación web.
\end{itemize}

\subsection{Transversales}
\begin{itemize}
\item Investigar y conocer los vectores de vigilancia usados habitualmente para
  monitorizar servicios de Internet.
\item Ampliar mis conocimientos sobre desarrollo web en general y las
  tecnologías de back-end en particular.
\item Adquirir soltura en el uso del lenguaje de programación Python en entornos
  web.
\item Obtener una base de conocimientos mínima sobre el desarrollo de
  aplicaciones sobre la plataforma móvil Android.
\item Utilizar un enfoque de análisis, diseño y codificación orientado
  a objetos, de una forma lo más clara y modular posible, para
  permitir ampliaciones y modificaciones sobre la aplicación por
  terceras personas.
\item Hacer uso de herramientas básicas en el desarrollo de software,
  como son los \textbf{Sistemas de Control de Versiones} para llevar
  un control realista del desarrollo del software, así como hacer de
  las veces de sistema de copias de seguridad.
\end{itemize}


\section{Alcance}
\textbf{SiteUp} se modela como una herramienta de monitorización de servicios de
internet accesible a través de la web. Los usuarios tendrán la posibilidad de
crear y gestionar una serie de \textit{chequeos} de diversos tipos sobre los
servicios web que elijan. La aplicación irá recopilando información relativa a
esos chequeos, e informará al usuario en caso de que las verificaciones que se
hayan dado de alta no coincidan con los resultados obtenidos.

Además, el usuario tendrá la posibilidad de recibir notificaciones de manera
instantánea a través del correo electrónico y de una aplicación para la
plataforma móvil Android. El servicio web estará totalmente adaptado para su uso
en dispositivos móviles.

\subsection{Limitaciones del proyecto}
Aunque cubre una gran parte de los puntos de vigilancia habituales, la
aplicación se limita a ofrecer chequeo de respuesta de ping, chequeo de puertos,
chequeo de registros DNS y chequeo de cabeceras y contenidos HTTP. 

La aplicación de Android no cuenta con ninguna funcionalidad para gestionar los
chequeos de un usuario, sino que sirve para recibir notificaciones instantáneas
provenientes de la aplicación web. Ésta, por otro lado, está completamente
adaptada para su uso a través de dispositivos móviles gracias al uso del
\textit{responsive web design}.

Idealmente los chequeos deberían hacerse simultáneamente desde diferentes
máquinas colocadas en diversos puntos geográficos, para así tener unos
resultados más fiables. La falta de infraestructuras y la finalidad didáctica
del proyecto han limitado la aplicación a una estructura monolítica en la que
los chequeos se hacen desde una sola máquina, la misma que sirve el servicio
web.

\subsection{Licencia}
El proyecto está publicado como software libre bajo la licencia
\ac{GPL} versión 3. El conjunto de bibliotecas y módulos utilizados
tienen las siguientes licencias:

\begin{itemize}

\item El framework de desarrollo web en el que se basa la aplicación es
  \textbf{Django}~\cite{django}, que utiliza la licencia
  \textit{\ac{BSD}}~\footnote{\url{https://github.com/django/django/blob/master/LICENSE}}.

\item El servidor web \textbf{nginx}~\cite{nginx} también utiliza la licencia \textit{BSD}.

\item Los siguientes paquetes de Python utilizan también la licencia \textit{BSD}:
  \begin{itemize}
  \item Celery
  \item Sqlparse
  \item iPython
  \item dnspython
  \item coverage
  \item django-rest-framework
  \item billiard
  \item anyjson
  \item Fabric
  \end{itemize}

\item Los siguientes paquetes de Python utilizan la licencia \textit{\ac{MIT}}:
  \begin{itemize}
  \item PyDot
  \item Gunicorn
  \item Requests
  \item django-extensions
  \item six
  \item sh
  \item pip
  \item virtualenvwrapper
  \item factory-boy
  \end{itemize}

\end{itemize}

\section{Estructura del documento}
El presente documento se rige según la siguiente estructura:

\begin{itemize}
\item \textbf{\nameref{chap:introduccion}}. Se exponen las motivaciones y
  objetivos detrás del proyecto \textbf{SiteUp}, así como información sobre las
  licencias de sus componentes, glosario y estructura del documento.
\item \textbf{\nameref{chap:calendario}}, donde se explica la planificación del
  proyecto, la división de sus etapas, la extensión de las etapas a lo largo del
  tiempo y los porcentajes de esfuerzo.
\item \textbf{\nameref{chap:fundamentos}}, que explica las labores de
  documentación y experimentación previas al desarrollo, que han servido para
  labrar una base de conocimientos que nos diera las suficientes garantías para
  afrontar el proyecto.
\item \textbf{\nameref{chap:analisis}}. Se detalla la fase de análisis del
  sistema, explicando los requisitos funcionales del sistema, los diferentes
  casos de uso, así como las principales operaciones con sus diagramas de
  secuencia y contratos.
\item \textbf{\nameref{chap:diseno}}. Seguido del análisis, se expone en detalle
  la etapa de diseño del sistema, con los diagramas de clases.
\item \textbf{\nameref{chap:implementacion}}. Una vez analizado el sistema y
  definido su diseño, en esta parte se detallan las decisiones de implementación
  más relevantes que tuvieron lugar durante el desarrollo del proyecto.
\item \textbf{\nameref{chap:pruebas}}. Listamos y describimos las pruebas que se
  han llevado a cabo sobre el proyecto para garantizar su fiabilidad y
  consistencia.
\item \textbf{\nameref{chap:conclusiones}}. Comentamos las conclusiones a las
  que se han llegado durante el transcurso y al término del proyecto.
\end{itemize}

Y los siguientes apéndices:
\begin{itemize}
\item \textbf{\nameref{chap:herramientas}}, donde detallamos qué hemos usado para la
  elaboración del proyecto.
\item \textbf{\nameref{chap:manual_instalacion}} del proyecto en sistemas nuevos.
\item \textbf{\nameref{chap:manual_usuario}}, donde se explica cómo usar la aplicación.
\end{itemize}

\subsection{Acrónimos}

\begin{acronym}
  \acro{IANA}{Internet Assigned Numbers Authority}
  \acro{IETF}{Internet Engineering Task Force}
  \acro{Sass}{Syntactically Awesome Style Sheets}
  \acro{AJAX}{Asynchronous JavaScript and XML}
  \acro{SQL}{Structured Query Language}
  \acro{ORM}{Object-relational mapping}
  \acro{AMQP}{Advanced Message Queuing Protocol}
  \acro{LOPD}{Ley Orgánica de Protección de Datos}
  \acro{WSGI}{Web Server Gateway Interface}
  \acro{PSFL}{Python Software Foundation License}
  \acro{ABI}{Application Binary Interface}
  \acro{ACL2s}{ACL2 Sedan}
  \acro{ACL2}{A Computational Logic for Applicative Common Lisp}
  \acro{ALSA}{Advanced Linux Sound Architecture}
  \acro{AMD}{Advanced Micro Devices}
  \acro{API}{Application Programming Interface}
  \acro{APL}{Apache Public License}
  \acro{ASCII}{American Standard Code for Information Interchange}
  \acro{BNF}{Backus-Naur Form}
  \acro{BOM}{Byte Order Mark}
  \acro{BSD}{Berkeley Software Distribution}
  \acro{C3}{Chrysler Comprehensive Compensation System}
  \acro{CAP}{Complex Arithmetic Processor}
  \acro{CASE}{Computer Assisted Software Engineering}
  \acro{CRUD}{Create-Read-Update-Delete}
  \acro{CD}{Compact Disc}
  \acro{CPAN}{Comprehensive Perl Archive Network}
  \acro{CPU}{Central Processing Unit}
  \acro{CSS}{Cascading Style Sheet(s)}
  \acro{CVS}{Concurrent Versions System}
  \acro{DDR}{Dance Dance Revolution}
  \acro{DNS}{Domain Name Server}
  \acro{DOM}{Document Object Model}
  \acro{DFT}{Discrete Fourier Transform}
  \acro{DSP}{Digital Signal Processing}
  \acro{DTD}{Do\-cu\-ment Type De\-fi\-ni\-tion}
  \acro{EAFP}{Easier to Ask for Forgiveness than Permission}
  \acro{FDL}{Free Documentation License}
  \acro{FFT}{Fast Fourier Transform}
  \acro{FSF}{Free Software Foundation}
  \acro{GCC}{GNU Compiler Collection}
  \acro{GCJ}{GNU Compiler for Java}
  \acro{GCL}{GNU Common Lisp}
  \acro{GIMP}{\acs{GNU} Image Manipulation Program}
  \acro{GML}{Generalized Markup Language}
  \acro{GNU}{GNU is Not Unix}
  \acro{GPL}{General Public License}
  \acro{GTK+}{\acs{GIMP} Toolkit}
  \acro{GUI}{Graphical User Interface}
  \acro{HTML}{Hyper Text Markup Language}
  \acro{HTTP}{Hyper Text Transfer Protocol}
  \acro{IBM}{International Business Machines}
  \acro{ICMP}{Internet Control Message Protocol}
  \acro{IDE}{Integrated Development Environment}
  \acro{IEEE}{Institute of Electrical and Electronics Engineers}
  \acro{IIS}{Internet Information Server}
  \acro{ISP}{Internet Service Provider}
  \acro{IP}{Internet Protocol}
  \acro{J2SE}{Java 2 Standard Edition}
  \acro{JACK}{JACK Audio Connection Kit}
  \acro{JAXP}{Java API for XML Parsing}
  \acro{JDBC}{Java DataBase Connectivity}
  \acro{JDK}{Java Development Kit}
  \acro{JIT}{Just In Time}
  \acro{JRE}{Java Runtime Environment}
  \acro{JSON}{JavaScript Object Notation}
  \acro{JSR}{Java Specification Request}
  \acro{JVM}{Java Virtual Machine}
  \acro{LGPL}{Lesser General Public License}
  \acro{LTS}{Long Time Support}
  \acro{MDI}{Multiple Document Interface}
  \acro{MIDI}{Musical Instrument Digital Interface}
  \acro{MIT}{Massachusetts Institute of Technology}
  \acro{MP3}{MPEG-2 Audio Layer III}
  \acro{MVC}{Model View Controller}
  \acro{MVP}{Model View Presenter}
  \acro{ODF}{Open Document Format}
  \acro{OSS}{Open Sound System}
  \acro{PAR}{Perl ARchive Toolkit}
  \acro{PCM}{Pulse-Code Modulation}
  \acro{PDF}{Portable Document Format}
  \acro{POD}{Plain Old Do\-cu\-men\-ta\-tion}
  \acro{POSIX}{Portable Operating System Interface, UniX based}
  \acro{Perl}{Practical Extraction and Report Language}
  \acro{RPM}{RPM Package Manager}
  \acro{RUP}{Rational Unified Process}
  \acro{RTP}{Real-time Transport Protocol}
  \acro{SAX}{Simple API for XML}
  \acro{SDI}{Single Document Interface}
  \acro{SDL}{Simple DirectMedia Layer}
  \acro{SGML}{Standard Generalized Markup Language}
  \acro{SSL}{Secure Sockets Layer}
  \acro{SVG}{Structured Vector Graphics}
  \acro{SWT}{Standard Widget Toolkit}
  \acro{TCK}{Technology Compatibility Kit}
  \acro{TDI}{Tabbed Document Interface}
  \acro{UML}{Unified Modelling Language}
  \acro{URI}{Uniform Resource Identifier}
  \acro{URL}{Uniform Resource Locator}
  \acro{UTF}{Universal Transformation Format}
  \acro{VPS}{Virtual Private Server}
  \acro{W3C}{World Wide Web Consortium}
  \acro{WWW}{World Wide Web}
  \acro{XHTML}{eXtensible Hyper Text Markup Language}
  \acro{XML}{eXtensible Markup Language}
  \acro{XP}{eXtreme Programming}
  \acro{XSL-FO}{eXtensible Stylesheet Language Formatting Objects}
  \acro{XSLT}{eXtensible Stylesheet Language Transformations}
  \acro{XSL}{eXtensible Stylesheet Language}
  \acro{YAML}{YAML Ain't a Markup Language}
  \acro{YAXML}{YAML XML binding}
\end{acronym}



%%% Local Variables:
%%% mode: latex
%%% TeX-master: "../memoria"
%%% End:


\chapter{Planificación}
\label{chap:calendario}
El proyecto no se ha desarrollado siguiendo un calendario estricto,
dado que era imposible cuantificar el tiempo que tomaría el adquirir
las bases teóricas necesarias para poder afrontarlo con garantías. Su
desarrollo se ha compaginado con los estudios del último curso de
Ingeniería Técnica en Informática de Sistemas y las labores como
becario en la Oficina de Software Libre y Conocimiento Abierto de la
Universidad de Cádiz~\cite{osluca}.

\section{Iteraciones}
Para la realización del proyecto se ha utilizado un modelo de desarrollo
iterativo incremental. En la redacción del presente documento se presentarán la
fase de investigación preliminar y las etapas de análisis, diseño,
implementación y pruebas del proyecto en su estado final. A continuación se
detallan cada una de las iteraciones por las que ha ido pasando el proyecto.

\subsection{Primera iteración: conocimientos preliminares}
Antes de poder comenzar con el análisis y diseño del propio proyecto,
era esencial adquirir una serie de conocimientos para poder afrontar
su desarrollo con todas las garantías. Durante esta iteración, se
llevaron a cabo labores de documentación y aprendizaje autodidacta con
las que se asentaron los conocimientos necesarios.

Además, durante este periodo también se barajaron las diferentes
posibilidades de implementación del sistema, así como las posibles
herramientas y bibliotecas de terceros que pudieran ser de ayuda.

\subsection{Segunda iteración: analizador básico}
Una vez adquiridos los conocimientos teóricos necesarios, y decididas
las técnicas y herramientas para llevar aquellos a la práctica, fue
obvia la necesidad de empezar por diseñar un analizador de notas
básico, que sería el corazón del programa. Del éxito del desarrollo
temprano del módulo que se encargaría del análisis de sonidos
dependería la viabilidad completa del proyecto.

\subsection{Tercera iteración: interfaz gráfica de usuario}
Con el módulo de análisis desarrollado, \textit{sólo} restaba
desarrollar el resto de la aplicación alrededor del mismo. En esta
tercera iteración se propusieron numerosos diseños para la interfaz
gráfica de usuario y, una vez decantados por uno de ellos, comenzó el
desarrollo de los elementos de la interfaz, haciendo énfasis en
conseguir un aspecto dinámico y jovial.

\subsection{Cuarta iteración: motor de lecciones}
Uno de los subproductos de la aplicación es el motor de lecciones, que
presenta una serie de unidades didácticas en formato multimedia,
compuestas de imágenes y textos, con conceptos sobre música. En esta
iteración se hizo un análisis de las posibilidades de este motor,
concluyendo con el diseño y desarrollo de un mecanismo muy sencillo de
ampliar y utilizar.

\subsection{Quinta iteración: motor de canciones}
La parte de mayor interactividad de la aplicación es el motor de
canciones, en el que el usuario tiene la posibilidad de tocar una
canción que aparece en pantalla, usando la flauta, mientras la
aplicación valora en tiempo real su interpretación. Durante la quinta
iteración se elaboró este sistema, encargado de listar y cargar las
diferentes canciones, y puntuar al usuario según cómo lo haga.

\section{Diagrama de Gantt}
Se ha diseñado un diagrama de Gantt para reflejar la distribución de las tareas
a lo largo del tiempo (figura~\ref{fig:gantt} en la página~\pageref{fig:gantt}).

\begin{figure}
  \centering
  \includegraphics[width=0.95\textwidth]{2_calendario/imagen_diagrama_gantt}
  \caption{Diagrama Gantt de iteraciones}
  \label{fig:gantt}
\end{figure}

%%% Local Variables: 
%%% mode: latex
%%% TeX-master: "../memoria"
%%% End: 


\chapter{Requisitos del sistema}
\label{chap:fundamentos}
En este capítulo se explicarán en detalle las instrucciones de uso de la
plataforma web.

\section{Introducción a SiteUp}

\textbf{SiteUp} es un servicio web creado para la monitorización de servicios de
Internet, especialmente pensado para dueños de sitios web y servidores que
deseen estar informados en caso de cualquier eventualidad.

El uso de internet en todos los ámbitos de la vida, desde las interacciones
sociales hasta la búsqueda de empleo, han dado lugar a la creación de servicios
accesibles a través de Internet y a nuevos modelos de negocio basados en
aquellos. Resulta por tanto imprescindible contar con sistemas para monitorizar
la disponibilidad de estos servicios y, en caso necesario, actuar de manera que
puedan subsanarse las causas de los problemas detectados.

SiteUp cubre precisamente estas necesidades, permitiendo a los usuarios:

\begin{itemize}
\item Crear chequeos mediante envío de paquetes ICMP (Ping), para revisar que un
  servidor está activo y conectado a la red.
\item Crear chequeos de registros DNS, de forma que el contenido de estos
  registros sea el correcto en todo momento.
\item Crear chequeos de puertos remotos, para vigilar que los servicios
  concretos dentro de un servidor son accesibles.
\item Crear chequeos mediante peticiones HTTP, para controlar que un servicio
  web responde de forma correcta y con el contenido adecuado.
\item Establecer las opciones de periodicidad de estos chequeos, de acuerdo a
  las necesidades particulares de cada caso.
\item Recibir notificaciones tanto a través de correo electrónico como mediante
  una aplicación móvil para el sistema operativo Android.
\end{itemize}

\section{Requisitos previos}

Para el uso correcto de la plataforma web SiteUp, el principal requisito es
contar con un dispositivo con acceso a internet que disponga de un navegador
reciente, como por ejemplo Google Chrome. Cualquier tipo de tablet, teléfono
móvil u ordenador personal podrá acceder a la plataforma web SiteUp, pero para
aprovechar mejor el potencial del sistema se recomienda usar una computadora
personal con las siguientes características:

\begin{itemize}
\item Navegador reciente. Se recomienda el uso de Google Chrome.
\item Pantalla con resolución mínima de 1280 píxeles en su lado longitudinal.
\item Conexión a internet con un caudal de descarga de al menos 512kb/s.
\end{itemize}

Si desea hacer uso de la aplicación móvil, será necesario que el usuario cuente
con un dispositivo con las siguientes características mínimas:

\begin{itemize}
\item Sistema operativo Android, versión mínima 2.3, versión recomendada 4.3.
\item Pantalla de al menos 3 pulgadas de diagonal y resolución mínima de 320x480
  píxeles.
\end{itemize}

\section{Uso de la plataforma web}

\subsection{Acceso}

A fecha de abril de 2014, la dirección de acceso a la plataforma web de SiteUp
es \url{http://siteup.josetomastocino.com}. Simplemente introduzca esa dirección
en su navegador y automáticamente será redirigido a la plataforma web, donde
verá la pantalla de bienvenida.

\subsection{Registro de usuario}

\begin{figure}[hbtp]
  \centering
  \includegraphics[width=0.3\textwidth]{apendice_manual_usuario/general_botones.png}
  \caption{Botones principales para el usuario anónimo}
  \label{fig:botones-1}
\end{figure}

Si es la primera vez que accede al sitio web deberá registrarse. Para ello, debe
pulsar en el botón \textbf{Sign up}, que se ve en la figura \ref{fig:botones-1}.

Tras ello, aparecerá el formulario de registro de usuario, que se ve en la
figura \ref{fig:registro}

\begin{figure}[hbtp]
  \centering
  \includegraphics[width=0.5\textwidth]{apendice_manual_usuario/formulario_registro.png}
  \caption{Formulario de registro de usuario}
  \label{fig:registro}
\end{figure}

En el formulario deberá introducir los siguientes datos:

\begin{itemize}
\item \textbf{Nombre de usuario}. Debe tener 30 caracteres como máximo, y solo
  pueden usarse letras, números y guiones.
\item \textbf{Dirección de correo electrónico}. Asegúrese de que es una
  dirección de correo correcta.
\item \textbf{Contraseña}.
\end{itemize}

Tras introducir los datos, deberá pulsar el botón \textbf{Sign up} del
formulario. Si el registro es correcto, el sistema dará de alta su usuario y
mostrará la pantalla \ref{fig:registro-correcto}. Si no, volverá a aparecer el
formulario con información sobre los errores en los datos.

\begin{figure}[hbtp]
  \centering
  \includegraphics[width=\textwidth]{apendice_manual_usuario/pantalla_register_completo.png}
  \caption{Mensaje indicando el registro correcto}
  \label{fig:registro-correcto}
\end{figure}

\subsection{Inicio de sesión}

Una vez que haya registrado su usuario, deberá pulsar en el botón \textbf{Log
  in}, como aparece en la imagen \ref{fig:botones-1}. Esto le llevará al
formulario de inicio de sesión que aparece en la figura \ref{fig:login}

\begin{figure}[hbtp]
  \centering
  \includegraphics[width=0.5\textwidth]{apendice_manual_usuario/general_formulario_login.png}
  \caption{Formulario de inicio de sesión}
  \label{fig:login}
\end{figure}

Deberá introducir su nombre de usuario y la contraseña que escribió al
registrarse, y posteriormente pulsar en el botón \textbf{Log in} del
formulario. Si los datos son correctos, será llevado al panel principal de la
web, donde podrá gestionar sus chequeos. Si no son correctos, volverá a aparecer
el formulario, con información sobre los errores encontrados.

\subsection{Recuperación de contraseña}

Si se registró como usuario pero no recuerda su contraseña es posible
recuperarla.

Para ello, vaya al formulario de inicio de sesión como ser ha explicado en el
paso anterior. En ese formulario, en la zona inferior, verá un enlace titulado
\textbf{Forgot password?} que deberá pulsar. Después de hacerlo se mostrará el
formulario de recuperación de contraseña, ilustrado en la figura
\ref{fig:formulario-reset}.

\begin{figure}[hbtp]
  \centering
  \includegraphics[width=0.5\textwidth]{apendice_manual_usuario/general_formulario_reset.png}
  \caption{Formulario de recuperación de contraseña}
  \label{fig:formulario-reset}
\end{figure}

En este formulario deberá introducir la dirección de correo electrónico que
utilizó al registrarse, y pulsar el botón \textbf{Send reset code}. Si lo hace
correctamente, aparecerá la pantalla de confirmación que se ve en la figura
\ref{fig:pantalla-reset-enviado}.

\begin{figure}[hbtp]
  \centering
  \includegraphics[width=0.5\textwidth]{apendice_manual_usuario/pantalla_reset_enviado.png}
  \caption{Pantalla de confirmación de envío de código de reinicio}
  \label{fig:pantalla-reset-enviado}
\end{figure}

El sistema habrá enviado un enlace a su dirección de correo electrónico. Deberá
pulsar el enlace que aparece en ese correo. Asegúrese de que recibe
correctamente el correo. Si no lo recibe, revise la carpeta de correo no
deseado. El enlace del correo le llevará a un formulario de reseteo de
contraseña como el que aparece en la figura \ref{fig:formulario-reset-2}.

\begin{figure}[hbtp]
  \centering
  \includegraphics[width=0.5\textwidth]{apendice_manual_usuario/general_formulario_reset_2.png}
  \caption{}
  \label{fig:formulario-reset-2}
\end{figure}

En este formulario deberá escribir la nueva contraseña de su cuenta dos veces,
para confirmar que la escribe correctamente. Tras ello, deberá pulsar el botón
\textbf{Set new password}, lo cual confirmará los cambios. 

Tras este proceso, podrá iniciar sesión como se ha descrito previamente, esta
vez utilizando la nueva contraseña.


\subsection{Pantalla de listado de chequeos}

\begin{figure}[hbtp]
  \centering
  \includegraphics[width=0.8\textwidth]{apendice_manual_usuario/pantalla-dashboard.png}
  \caption{Panel de control}
  \label{fig:dashboard}
\end{figure}

\begin{figure}[hbtp]
  \centering
  \includegraphics[width=0.4\textwidth]{apendice_manual_usuario/general_botones_2.png}
  \caption{Botones principales para el usuario logueado}
  \label{fig:botones-2}
\end{figure}

En la figura \ref{fig:dashboard} se puede ver el aspecto del panel de control de
la plataforma web SiteUp. Desde este panel podrá hacer una gran cantidad de
acciones, desde crear grupos de chequeos hasta editar sus detalles de usuario.

Es posible acceder a esta pantalla desde cualquier sitio de la web si está
logueado, pulsando en el botón \textbf{Dashboard} de la cabecera de la web, tal
y como se ve en la figura \ref{fig:botones-2}.

\subsection{Cerrar sesión}

Si ha iniciado sesión y desea cerrarla, debe dirigirse a los botones
principales, situados en el lado derecho de la cabecera de la web, tal y como
aparecen en la figura \ref{fig:botones-2}. Deberá pulsar el botón \textbf{Log
  out}, tras lo cual se cerrara su sesión y será redirigido a la pantalla
inicial de SiteUp.

\subsection{Edición de datos de usuario}

Si desea editar alguno de los datos de usuario que ingresó cuando hizo el
registro, deberá primero pulsar el botón \textbf{Your profile}, situado en la
cabecera de la web. Es el segundo de los botones que aparecen, tal y como se ve
en la figura \ref{fig:botones-2}.

Una vez hecho esto, aparecerá el formulario de edición de usuario, ilustrado en
la figura \ref{fig:formulario-usuario}, en el que verá sus actuales datos de
usuario. 

\begin{figure}[hbtp]
  \centering
  \includegraphics[width=0.8\textwidth]{apendice_manual_usuario/pantalla_perfil_usuario.png}
  \caption{}
  \label{fig:formulario-usuario}
\end{figure}

En este formulario podrá editar:

\begin{itemize}
\item Su \textbf{nombre} de usuario.
\item Su \textbf{dirección} de correo electrónico.
\item Si desea o no recibir \textbf{resúmenes diarios} sobre el estado de sus chequeos.
\end{itemize}

Una vez editados los datos que desee deberá pulsar el botón \textbf{Update
  details}. También podrá editar su contraseña pulsando el botón \textbf{Change
  password}, situado al pie del formulario a la izquierda.


\subsection{Creación de grupo de chequeos}

En SiteUp, los chequeos se organizan en \textbf{grupos de chequeos}, lo que
permite trabajar con ellos más fácilmente a la hora de activarlos, desactivarlos
o borrarlos. 

Para crear un grupo de chequeos deberá pulsar el botón \textbf{Add group},
situado en la zona superior derecha del panel principal. El botón puede verse en
la figura \ref{fig:crear-grupo}.

\begin{figure}[hbtp]
  \centering
  \includegraphics[width=0.3\textwidth]{apendice_manual_usuario/grupo_boton_nuevo.png}
  \caption{Botón para crear nuevo grupo de chequeos}
  \label{fig:crear-grupo}
\end{figure}

Tras pulsar el botón aparecerá el formulario de creación de nuevo grupo. En él
deberá introducir los datos del nuevo grupo. En particular:

\begin{itemize}
\item El \textbf{título} que identifica al grupo.
\end{itemize}

Tras ello, deberá pulsar el botón \textbf{Create group}, y el grupo será creado

\begin{figure}[H]
  \centering
  \includegraphics[width=0.8\textwidth]{apendice_manual_usuario/grupo_formulario_crear.png}
  \caption{Formulario de creación de grupo de chequeos}
  \label{fig:}
\end{figure}


\subsection{Operaciones con grupos de chequeos}

En la figura~\ref{fig:operaciones-grupo} se puede ver el menú de operaciones
disponibles para un grupo de chequeos. Este menú aparece en el panel principal,
al lado derecho junto al título de cada grupo de chequeos.

\begin{figure}[hbtp]
  \centering
  \includegraphics[width=0.5\textwidth]{apendice_manual_usuario/grupo_submenu_acciones.png}
  \caption{}
  \label{fig:operaciones-grupo}
\end{figure}

\subsubsection{Edición de grupo de chequeos}

En el menú de operaciones del grupo de chequeos, pulse el botón
\textbf{Edit}. Aparecerá un formulario similar al de creación de grupo, en el
que podrá editar el título del grupo de chequeos.

\subsubsection{Borrado de grupo de chequeos}

Para borrar un grupo de chequeos, pulse el botón \textbf{Delete} en el menú de
operaciones del grupo. Aparecerá un formulario como el de la figura
\ref{fig:formulario-borrado-grupo}, pidiéndole confirmación del borrado. 

\begin{figure}[hbtp]
  \centering
  \includegraphics[width=0.5\textwidth]{apendice_manual_usuario/formulario_borrado_grupo}
  \caption{Formulario de confirmación de borrado}
  \label{fig:formulario-borrado-grupo}
\end{figure}

Si pulsa en el botón \textbf{Go back}, volverá al panel principal, cancelando el
borrado. Si pulsa el botón \textbf{Delete}, confirmará la operación y el sistema
borrará el grupo de chequeos, así como todos los chequeos que contenga.

\subsubsection{Activación y desactivación de grupo de chequeos}

Es posible activar y desactivar de forma grupal todos los chequeos de un
grupo. Para ello, deberá dirigirse al menú de operaciones del grupo y pulsar en
\textbf{Enable all checks} o en \textbf{Disable all checks}, según sea su decisión.

\subsection{Creación de chequeo}

\begin{figure}[hbtp]
  \centering
  \includegraphics[width=0.8\textwidth]{apendice_manual_usuario/pantalla_seleccion_tipo.png}
  \caption{Pantalla de selección del tipo de chequeo a crear}
  \label{fig:seleccionar-tipo-check}
\end{figure}

Como se ha comentado previamente, el elemento principal de SiteUp es el
\textbf{chequeo}. El usuario da de alta chequeos dentro de un grupo, y estos
chequeos se lanzan periódicamente, generando datos.

Para crear un chequeo dentro de un grupo deberá ir al menú de operaciones del
grupo, que aparece en la figura \ref{fig:operaciones-grupo} y se sitúa a la
derecha del título de un grupo de chequeos. Dentro del menú, deberá pulsar la
opción \textbf{Add check}.



Tras esto aparecerá una pantalla en la que deberá seleccionar el tipo de chequeo
que desea crear. La pantalla es similar a la de la figura
\ref{fig:seleccionar-tipo-check}. Como se explicó en la introducción de este
manual, es posible crear cuatro tipo de chequeos: chequeos de ping, chequeos de
puerto, chequeos de registros DNS y chequeos por HTTP.

Si pulsa en el botón del tipo que desea crear, aparecerá un formulario con
campos para rellenar.

\subsubsection{Creación de un chequeo por ping}

\begin{figure}[hbtp]
  \centering
  \includegraphics[width=0.8\textwidth]{apendice_manual_usuario/pantalla_crear_ping.png}
  \caption{Formulario de creación de chequeo por ping}
  \label{fig:crear-ping}
\end{figure}

El formulario para crear un chequeo de tipo ping aparece en la
figura~\ref{fig:crear-ping}. Los campos que deberá rellenar son los siguientes:

\begin{itemize}
\item \textbf{Target}: es el objetivo a chequear. Debe ser un \textit{hostname} o una IP.
\item \textbf{Título} del chequeo.
\item \textbf{Descripción} opcional del chequeo.
\item \textbf{Check interval}: intervalo entre chequeos, definido en minutos. El
  valor mínimo es de 1 minuto.
\item \textbf{Notify email}: indica si los cambios de estado en este chequeo
  deberán ser notificados por correo electrónico.
\item \textbf{Notify Android}: indica si los cambios de estado en este chequeo
  deberán ser notificados mediante la aplicación Android.
\item \textbf{Should check timeout}: indica si hay que verificar que el ping
  enviado vuelve dentro de un tiempo máximo determinado.
\item \textbf{Timeout value}: si la anterior opción está activa, este campo
  define el número de milisegundos que puede tardar en responder el servidor al
  ping, como máximo.
\end{itemize}

\subsubsection{Creación de un chequeo de puerto}

\begin{figure}[hbtp]
  \centering
  \includegraphics[width=0.8\textwidth]{apendice_manual_usuario/pantalla_crear_port.png}
  \caption{Formulario de creación de chequeo de puerto}
  \label{fig:crear-port}
\end{figure}

El formulario para crear un chequeo de puertos aparece en la
figura~\ref{fig:crear-port}. Los campos que deberá rellenar son los siguientes:

\begin{itemize}
\item \textbf{Título} del chequeo.
\item \textbf{Descripción} opcional del chequeo.
\item \textbf{Target}: es el objetivo a chequear. Debe ser un \textit{hostname} o una IP.
\item \textbf{Target port}: es el puerto a chequear. Debe ser un entero.
\item \textbf{Check for string}: opcionalmente, se puede indicar una cadena que
  debe estar presente en la respuesta al conectar con el puerto.
\item \textbf{Check interval}: intervalo entre chequeos, definido en minutos. El
  valor mínimo es de 1 minuto.
\item \textbf{Notify email}: indica si los cambios de estado en este chequeo
  deberán ser notificados por correo electrónico.
\item \textbf{Notify Android}: indica si los cambios de estado en este chequeo
  deberán ser notificados mediante la aplicación Android.
\end{itemize}

\subsubsection{Creación de un chequeo de DNS}

\begin{figure}[hbtp]
  \centering
  \includegraphics[width=0.8\textwidth]{apendice_manual_usuario/pantalla_crear_dns.png}
  \caption{Formulario de creación de chequeo de registro DNS}
  \label{fig:crear-dns}
\end{figure}

El formulario para crear un chequeo de registro DNS aparece en la figura
\ref{fig:crear-dns}. Los campos que deberá rellenar son los siguientes:

\begin{itemize}
\item \textbf{Título} del chequeo.
\item \textbf{Descripción} opcional del chequeo.
\item \textbf{Target}: es el objetivo a chequear. Debe ser un \textit{hostname}.
\item \textbf{Record type}: indica el tipo de registro a chequear. Puede ser A, AAAA, CNAME, TXT y MX.
\item \textbf{Resolved address}: indica el contenido que debe tener el registro.
\item \textbf{Check interval}: intervalo entre chequeos, definido en minutos. El
  valor mínimo es de 1 minuto.
\item \textbf{Notify email}: indica si los cambios de estado en este chequeo
  deberán ser notificados por correo electrónico.
\item \textbf{Notify Android}: indica si los cambios de estado en este chequeo
  deberán ser notificados mediante la aplicación Android.
\end{itemize}

\subsubsection{Creación de un chequeo por HTTP}

\begin{figure}[hbtp]
  \centering
  \includegraphics[width=0.8\textwidth]{apendice_manual_usuario/pantalla_crear_http.png}
  \caption{Formulario de creación de chequeo por HTTP}
  \label{fig:crear-http}
\end{figure}

El formulario para crear un chequeo por HTTP aparece en la figura
\ref{fig:crear-http}. Los campos que deberá rellenar son los siguientes:

\begin{itemize}
\item \textbf{Título} del chequeo.
\item \textbf{Descripción} opcional del chequeo.
\item \textbf{Target}: es el objetivo a chequear. Debe ser una URL.
\item \textbf{Status code}: es el código de estado que deberá devolver la petición..
\item \textbf{Check for string}: opcionalmente, se puede indicar una cadena que
  debe estar presente en la respuesta a la petición.
\item \textbf{Check interval}: intervalo entre chequeos, definido en minutos. El
  valor mínimo es de 1 minuto.
\item \textbf{Notify email}: indica si los cambios de estado en este chequeo
  deberán ser notificados por correo electrónico.
\item \textbf{Notify Android}: indica si los cambios de estado en este chequeo
  deberán ser notificados mediante la aplicación Android.
\end{itemize}


\subsection{Operaciones con chequeos}

\begin{figure}[hbtp]
  \centering
  \includegraphics[width=\textwidth]{apendice_manual_usuario/chequeo_valido.png}
  \caption{}
  \label{fig:chequeo}
\end{figure}

Una vez creado un chequeo, aparecerá en el panel de control con el aspecto de la
figura~\ref{fig:chequeo}. Como se aprecia en la figura, se muestra mediante un
icono el estado del chequeo, los datos (título, tipo, descripción) y una gráfica
que ilustra el estado del chequeo en las últimas 24 horas.

\begin{figure}[hbtp]
  \centering
  \includegraphics[width=\textwidth]{apendice_manual_usuario/chequeo_invalido.png}
  \caption{}
  \label{fig:chequeo-invalido}
\end{figure}

Si un chequeo detecta un error, cambia de apariencia, pasando a tener el aspecto
que ofrece la figura \ref{fig:chequeo-invalido}. Además, según se haya
configurado en las opciones del chequeo, se envía una notificación al usuario
informando del cambio de estado del chequeo.

Cuando se ve en la lista general de chequeos, cada chequeo tiene una serie de
botones que permiten llevar a cabo operaciones diversas, que se detallan a
continuación.

\subsubsection{Activación y desactivación de un chequeo}

Para desactivar un chequeo, tiene que pulsar el botón \textbf{Disable} de la
botonera, que solo aparecerá si el chequeo está activado. De igual modo, si el
chequeo está desactivado, el botón que aparecerá tendrá la etiqueta
\textbf{Enable}, que podrá pulsar para activarlo.

\subsubsection{Borrado de un chequeo}

En la botonera se sitúa el botón \textbf{Delete}, que podrá pulsar para borrar
el chequeo. Aparecerá un formulario pidiéndole confirmación. Si confirma la
acción, se procederá al borrado del chequeo y todos los datos dependientes de
éste.

\subsubsection{Edición de un chequeo}

Si desea editar un chequeo debe pulsar el botón \textbf{Edit}, situado en la
botonera del chequeo. Aparecerá un formulario similar al de creación de chequeo,
pero con los datos actuales ya introducidos. Puede modificar los datos que desee
y confirmar los cambios pulsando el botón \textbf{Update check} del formulario.

\subsubsection{Vista de detalle del chequeo}

\begin{figure}[hbtp]
  \centering
  \includegraphics[width=\textwidth]{apendice_manual_usuario/pantalla_detalle_chequeo.png}
  \caption{}
  \label{fig:detalle-chequeo}
\end{figure}

En la botonera del chequeo, en el panel de control, verá un botón
\textbf{Details}. Si lo pulsa será llevado a la pantalla de vista de detalle del
chequeo, que se ilustra en la figura \ref{fig:detalle-chequeo}. Esta pantalla
ofrece mucha información sobre el chequeo. En particular:

\begin{itemize}
\item En la parte superior se sitúa la información general:
  \begin{itemize}
  \item Icono con el estado del chequeo.
  \item Título.
  \item Tipo de chequeo.
  \item Descripción.
  \item Botonera de acciones.
  \end{itemize}

\item En la parte media de la pantalla se sitúa un bloque con pestañas. Mediante
  estas pestañas es posible indicar sobre qué periodo de tiempo queremos
  visualizar los datos: últimas 24 horas o última semana. La información
  disponible es la siguiente:

  \begin{itemize}
  \item Porcentaje de \textit{uptime} del chequeo.
  \item \textit{(Solo chequeos tipo ping)} Tiempo de respuesta mínimo, máximo y promedio.
  \item Gráfica con el estado del chequeo.
  \end{itemize}

\item En la parte inferior de la pantalla se sitúa el panel de eventos del
  chequeo. Éste permite ver los últimos cambios de estado del chequeo:

  \begin{itemize}
  \item Tipo de cambio de estado (\textit{Up}, \textit{Down}).
  \item Duración del estado.
  \item Hora de comienzo del estado.
  \item Hora de fin del estado. Si es el estado en el que el chequeo se
    encuentra actualmente, no aparecerá ninguna fecha.
  \item Extra: muestra información adicional para estados de error.
  \end{itemize}
\end{itemize}

\begin{figure}[hbtp]
  \centering
  \includegraphics[width=\textwidth]{apendice_manual_usuario/pantalla_grafica_semanal.png}
  \caption{Detalle de la gráfica para un intervalo de tiempo de una semana}
  \label{fig:grafica-semanal}
\end{figure} 

\subsection{Recepción de notificaciones por correo electrónico}

\subsubsection{Notificaciones de cambio de estado}

Cuando el sistema detecte un cambio de estado en un chequeo, y esté activada la
notificación por correo electrónico, se enviará un email como el que aparece en
la figura \ref{fig:email-estado}.

\begin{figure}[hbtp]
  \centering
  \includegraphics[width=0.7\textwidth]{apendice_manual_usuario/email_estado.png}
  \caption{Email de notificación de cambio de estado}
  \label{fig:email-estado}
\end{figure}

En ese email se le indicará el título del chequeo, el nuevo estado, la fecha y
hora en la que ha pasado al nuevo estado, y un enlace directo al panel de
detalle del chequeo en SiteUp, de forma que pueda obtener más información.

\subsubsection{Notificaciones de resumen diarias}

Si ha activado el envío de resúmenes diarios sobre el estado de sus chequeos, el
sistema enviará todos los días un correo electrónico con información. El aspecto
del correo que recibirá se puede ver en la figura \ref{fig:email-resumen}

\begin{figure}[hbtp]
  \centering
  \includegraphics[width=0.7\textwidth]{apendice_manual_usuario/email_resumen.png}
  \caption{Email diario de resumen}
  \label{fig:email-resumen}
\end{figure}


\section{Uso de la aplicación móvil}

La aplicación móvil \textbf{SiteUp Client} es una aplicación para para el
sistema operativo Android que permite a sus usuarios revisar los chequeos que
tiene dados de alta en la plataforma web SiteUp, así como recibir de manera
instantánea notificaciones \textit{push} sobre el estado de sus chequeos.

\subsection{Instalación}

La aplicación móvil se instala copiando y ejecutando el fichero APK en el
dispositivo. En un futuro se podrá encontrar en la \textit{Play Store} de
Google.

\subsection{Lanzamiento}

En la figura~\ref{fig:manual-android-launcher} se puede observar el aspecto de
la aplicación \textbf{SiteUp Client} en el lanzador de Android. Para ejecutarla,
simplemente pulse sobre el icono de la aplicación.

\begin{figure}[htbp]
  \centering
  \includegraphics[width=0.4\textwidth]{5_diseno/android-1}
  \caption{Icono de SiteUp en el lanzador}
  \label{fig:manual-android-launcher}
\end{figure}

\subsection{Carga de la aplicación}

Tras lanzar la aplicación, el sistema cargará los datos necesarios y se
comunicará con los servidores de Google para obtener los datos necesarios. La
apariencia de la aplicación durante este proceso se puede ver en la figura
\ref{fig:manual-android-login}.

\begin{figure}[htbp]
  \centering
  \includegraphics[width=0.4\textwidth]{5_diseno/android-2}
  \caption{Pantalla de carga de la aplicación móvil}
  \label{fig:manual-android-loading}
\end{figure}

\subsection{Inicio de sesión}

Tras la carga, la primera vez que lance la aplicación será necesario que inicie
sesión con las credenciales que utilizó en la plataforma web. Para ello, tendrá
que introducir su nombre de usuario y contraseña en el formulario que aparece en
la figura~\ref{fig:manual-android-loading}, y pulsar el botón \textbf{Sign in}.

\begin{figure}[htbp]
  \centering
  \includegraphics[width=0.4\textwidth]{5_diseno/android-3}
  \caption{Pantalla de inicio de sesión de la aplicación móvil}
  \label{fig:manual-android-login}
\end{figure}

La aplicación móvil se comunicará con los servidores centrales de SiteUp,
validará sus credenciales de usuario y, si son correctas, le dará acceso al
sistema. Si las credenciales de usuario no son correctas, el formulario mostrará
un error.

\subsection{Revisión de chequeos}

Si ha iniciado sesión correctamente, tendrá la opción de revisar la lista de
chequeos que tiene dados de alta en la plataforma web \textbf{SiteUp}. La
apariencia de la lista de chequeos es igual a la de la
figura~\ref{fig:manual-android-dashboard}.

\begin{figure}[hbtp]
  \centering
  \includegraphics[width=0.4\textwidth]{5_diseno/android-4}
  \caption{Pantalla de chequeos de la aplicación móvil}
  \label{fig:manual-android-dashboard}
\end{figure}

Si pulsa en cualquiera de los chequeos, será redirigido a la plataforma web, que
podrá navegar sin problema gracias a que se encuentra preparada para su uso en
dispositivos móviles.

\subsection{Recepción de notificaciones}

La utilidad principal de la aplicación es la posibilidad de recibir
notificaciones en cualquier momento provenientes del servidor -- conocidas como
notificaciones \textit{push}. 

El procedimiento es transparente al usuario. Simplemente lance la aplicación e
inice sesión. Tras hacer esto, la aplicación estará preparada para recibir
notificaciones desde el servidor.

Cuando se reciba una notificación, la aplicación mostará un mensaje en el área
de notificaciones del móvil, con información resumida. Puede ver el aspecto de
la notificación en la figura \ref{fig:manual-android-notification}. Si pulsa en
la notificación será llevado a los detalles del chequeo para obtener más
detalles sobre el origen del aviso.

\begin{figure}[hbtp]
  \centering
  \includegraphics[width=0.4\textwidth]{5_diseno/android-5}
  \caption{Aspecto de las notificaciones}
  \label{fig:manual-android-notification}
\end{figure}



%%% Local Variables:
%%% mode: latex
%%% TeX-master: "../memoria"
%%% End:


\chapter{Análisis}
\label{chap:analisis}

\section{Metodología}
\textbf{oFlute} ha seguido una metodología de desarrollo ágil en la que,
mediante fases de desarrollo rápidas y ligeras, se intenta evitar los formales
caminos de las metodologías tradicionales, enfocándose en las personas y
obteniendo así resultados en etapas más tempranas del desarrollo.

En las sucesivas secciones y capítulos se detallará el proceso de análisis y
posteriores fases del proyecto en la última iteración del proyecto. Así se
consigue una documentación más concisa y cercana al producto final.

\section{Especificación de requisitos del sistema}

\subsection{Requisitos de interfaces externas}
En esta sección describiremos los requisitos que deben cumplir las interfaces
con el hardware, el software y el usuario.

En cuanto a la comunicación con el subsistema gráfico y de E/S, utilizaremos la
biblioteca Gosu~\cite{gosu}, un proyecto de software libre que proporciona un
framework de desarrollo de videojuegos 2D, multiplataforma y muy sencillo de
usar. Para el acceso al subsistema de audio, tal y como se ha comentado en la
sección anterior, optamos por utilizar la API simple de PulseAudio.

\textbf{oFlute} dispondrá de una resolución fija de 800 por 600 píxeles,
requisito fácilmente alcanzable en cualquier ordenador actual. Al tratar con un
público objetivo joven, los gráficos y la interactividad deberán ser sencillos y
fáciles de interpretar. Así, se ha trabajado en limitar la interacción del
usuario con la aplicación al uso del ratón y, obviamente, del instrumento
musical, en este caso la flauta dulce. La navegación resultante de este
planteamiento queda reflejada en el siguiente diagrama:

\begin{figure}[h!]
  \centering
  \includegraphics[width=0.9\textwidth]{4_analisis/imagen_diagrama_de_flujo}
  \caption{Diagrama de flujo de las pantallas de oFlute}
\end{figure}

\pagebreak

Inicialmente, deberán aparecer unas pantallas de crédito con información sobre
el desarrollador y sobre el propio videojuego. Tras las mismas, que deberá ser
posible omitir, habrá de aparecer el \textbf{menú principal}, con las cinco
opciones posibles.

\begin{figure}[h!]
  \centering
  \includegraphics[width=0.6\textwidth]{4_analisis/imagen_mockup_menu_principal}
  \caption{Maqueta del menú principal}
\end{figure}

\pagebreak

Las opciones que se incluirán son:
\begin{itemize}
\item \textbf{Analizador de notas}: comprobar las notas que tocamos sobre un
  pentagrama.
\item \textbf{Canciones}: sección principal del juego, en el que aparecerán las
  canciones a tocar.
\item \textbf{Lecciones}: sección de lecciones de aprendizaje.
\item \textbf{Calibrar micrófono}, para ajustarse al nivel de ruido ambiental.
\item \textbf{Salir} al sistema operativo.
\end{itemize}

La siguiente pantalla a modelar será el \textbf{analizador de
  notas}. Simplemente mostrará el logotipo del videojuego a un lado, y un
pentagrama al otro, que se actualizará con la nota detectada por el
micrófono. También contendrá un botón \textit{volver} para ir al menú principal.

\begin{figure}[h!]
  \centering
  \includegraphics[width=0.6\textwidth]{4_analisis/imagen_mockup_analizador}
  \caption{Maqueta de la sección \textit{analizador de notas}}
\end{figure}

La segunda sección a la que se podrá ir desde el menú principal será la de
\textbf{canciones}. Inicialmente, la primera pantalla será la de
\textbf{selección de canción}, que contendrá el logotipo del juego, un botón
para volver al menú principal, y un menú dinámico de canciones que nos permitirá
elegir el tema a interpretar.

\begin{figure}[h!]
  \centering
  \includegraphics[width=0.6\textwidth]{4_analisis/imagen_mockup_seleccionar_cancion}
  \caption{Maqueta del menú de selección de canción}
\end{figure}

Una vez seleccionada la canción, pasaremos a la zona de \textbf{interpretación
  de canción}. Contendrá un pentagrama que ocupará todo el ancho de la pantalla,
con una línea que indicará la zona donde empezar a tocar las notas que
aparezcan. Además, en la parte superior habrá un indicador con la puntuación
obtenida y, abajo, una barra de progreso que nos indicará cuánto queda de
canción.

\begin{figure}[h!]
  \centering
  \includegraphics[width=0.6\textwidth]{4_analisis/imagen_mockup_reproducir_cancion}
  \caption{Maqueta de la pantalla de interpretación de canción}
\end{figure}

\pagebreak

Al completar la interpretación de la canción, aparecerá la \textbf{sección de
  resultados}. Contendrá el logotipo del juego, el título y subtítulo de la
canción, y un cuadro con información sobre nuestra interpretación, representada
en forma de porcentaje de aciertos. Además, en la zona inferior aparecerá un
mensaje de ánimo dependiendo del resultado obtenido.

\begin{figure}[h!]
  \centering
  \includegraphics[width=0.6\textwidth]{4_analisis/imagen_mockup_puntuaciones}
  \caption{Maqueta de la pantalla de puntuaciones}
\end{figure}

La pantalla de \textbf{selección de lecciones}, a la que se llega desde el menú
principal, contendrá el título, una imagen decorativa, y varios botones para
navegar entre las lecciones cargadas en el sistema. Se mostrará el título y la
descripción de cada lección, así como un botón para comenzar.

\begin{figure}[h!]
  \centering
  \includegraphics[width=0.6\textwidth]{4_analisis/imagen_mockup_seleccionar_leccion}
  \caption{Maqueta del menú de selección de lecciones}
\end{figure}

Una vez elegida una lección, pasaremos a la pantalla de reproducción de
lecciones. Dada la naturaleza \textbf{dinámica} de esta sección, cada lección
podrá tener una apariencia y elementos distintos. El único elemento común entre
todas las lecciones será el botón de \textbf{volver al menú}.

\subsection{Requisitos funcionales}

\textbf{oFlute} se basa en los siguientes requisitos funcionales:
\begin{itemize}
\item Poder terminar la aplicación pulsando el botón de cierre en cualquier
  instante.
\item Comprobar la correcta interpretación de notas individuales mediante el
  analizador de notas.
\item Calibrar el micrófono de forma que el sistema se pueda adaptar al ruido
  ambiental del entorno.
\item Navegar por toda la aplicación de forma sencilla utilizando solo el ratón.
\item Elegir entre varias canciones a interpretar, cada una con su título y
  subtítulo informativos.
\item Interpretar las canciones mediante el uso de la flauta, siguiendo el
  pentagrama en pantalla.
\item Elegir entre bastantes lecciones informativas, poder ejecutarlas y
  seguirlas.
\item Capacidad de añadir nuevas lecciones y canciones de forma sencilla.
\end{itemize}

\subsection{Requisitos de rendimiento}

La aplicación \textbf{oFlute} precisa de unos requisitos bastante básicos,
que en su mayor parte se reducen a cuatro puntos principales:
\begin{itemize}
\item Posesión de una tarjeta de sonido o subsistema de audio similar con un
  micrófono, para poder captar el sonido de la flauta.
\item Pantalla con una resolución de, al menos, 800 por 600 píxeles.
\item Sistema gráfico compatible con OpenGL.
\item Dispositivo apuntador, como un ratón.
\end{itemize}

La práctica totalidad de los ordenadores personales de la actualidad cumplen los
citados requisitos.

\subsection{Requisitos del sistema software}
El sistema de software deberá cumpliar los requisitos siguientes:
\begin{itemize}
\item Deberá funcionar en cualquier sistema \textbf{GNU/Linux} con los
  requisitos anteriormente indicados.
\item Deberá limitarse el número de dependencias, así como facilitar al máximo
  la instalación de las que resultasen imprescindibles.
\item El uso del teclado quedará en segundo plano, haciendo posible utilizar la
  aplicación completamente con el ratón.
\item Al tratarse de un público objetivo juvenil, la aplicación deberá ser
  dinámica, intuitiva y fácil de usar, y la apariencia debe ser agradable.
\item Se evitará el uso de constantes y recursos dentro del código de la
  aplicación, utilizando como alternativa ficheros para representar las
  lecciones y las canciones.
\end{itemize}

\section{Modelo de casos de uso}

A la hora de modelar los casos de uso del sistema, hemos optado por utilizar
notación \textit{UML}, siguiendo los siguientes pasos:
\begin{itemize}
\item Identificación de los usuarios del sistema y sus roles.
\item Para cada rol, determinar las formas de interactuar con el sistema.
\item Creación de casos de uso para los objetivos que debe cumplir la aplicación.
\item Modularización de los casos de usos mediante la implementación de
  relaciones de inclusión o extensión.
\end{itemize}

\subsection{Diagrama de casos de uso}

\begin{figure}[h!]
  \centering
  \includegraphics[width=0.81\textwidth]{4_analisis/imagen_diagrama_de_casos_de_uso}
  \caption{Diagrama de casos de uso}
\end{figure}


\subsection{Descripción de los casos de uso}

\subsubsection{Caso de uso: inicio del juego}
\begin{description}
\item [Descripción] Se muestran los créditos del juego, la pantalla de
  presentación, y finalmente el menú principal, desde donde se accederá al
  resto de secciones del juego.
\item [Actores] \jugador.
\item [Precondiciones] Ninguna.
\item [Postcondiciones] Ninguna.
\item [Escenario principal] $\quad$
  \begin{enumerate}
  \item El \jugador\ inicia la aplicación.
  \item El \sistema\ inicializa el subsistema gráfico.
  \item El \sistema\ muestra la pantalla de créditos y la pantalla de
    presentación de la aplicación.
  \item El \sistema\ muestra el menú principal en la pantalla.
  \item El \jugador\ selecciona la opción \textit{Canciones}.
  \item El \sistema\ accede a la pantalla de \textit{Selección de canción}.
  \end{enumerate}
\item[Extensiones --- flujo alternativo] $\quad$
  \begin{description}
  \item [*a] El \jugador\ cierra la ventana.
    \begin{enumerate}
    \item El \sistema\ libera los recursos y sale de la aplicación.
    \end{enumerate}
  \item [4a] El \jugador\ selecciona la opción \textit{Analizador de Notas}.
    \begin{enumerate}
    \item El \sistema\ accede a la pantalla del analizador de notas.
    \end{enumerate}

  \item[4b] El \jugador\ selecciona la opción \textit{Lecciones}.
    \begin{enumerate}
    \item El \sistema\ accede a la pantalla de \textit{Selección de lecciones}.
    \end{enumerate}
  \item[4c] El \jugador\ selecciona la opción \textit{Calibrar micrófono}.
    \begin{enumerate}
    \item El \sistema\ accede a la pantalla de calibración de micrófono.
    \end{enumerate}
  \item [4d] El \jugador\ selecciona la opción Salir.
    \begin{enumerate}
    \item El \sistema\ libera los recursos y sale de la aplicación.\\
    \end{enumerate}
  \end{description}  
\end{description}


\subsubsection{Caso de uso: selección de canción}

\begin{description}
\item [Descripción] Al \jugador\ se le muestra una lista de las canciones
  detectadas, y éste debe elegir entre ellas la que desea interpretar, o volver
  al menú princial.
\item [Actores] \jugador.
\item [Precondiciones] Ninguna.
\item [Postcondiciones] Una canción queda seleccionada.
\item [Escenario principal] $\quad$
  \begin{enumerate}
  \item El \jugador\ accede, desde el menú principal, al panel de selección de canciones.
  \item El \sistema\ busca las canciones dadas de alta en el juego y muestra un menú con las mismas.
  \item El \jugador\ navega entre las canciones listadas y selecciona una de ellas, pulsando finalmente el botón \textit{Ok}.
  \item El \sistema\ carga la canción y pasa a la pantalla de interpretación de canciones.
  \end{enumerate}
\item[Extensiones --- flujo alternativo] $\quad$
  \begin{description}
  \item [*a] El \jugador\ cierra la ventana.
    \begin{enumerate}
    \item El \sistema\ libera los recursos y sale de la aplicación.
    \end{enumerate}
  \item [3a] El \jugador\ selecciona la opción \textit{Volver}.
    \begin{enumerate}
    \item El \sistema\ muestra la animación de cierre y vuelve al menú principal.
    \end{enumerate}
  \end{description}  
\end{description}


\subsubsection{Caso de uso: interpretación de canción}

\begin{description}
\item [Descripción] Tras haber elegido la canción a interpretar, se muestra una
  partitura con las notas que el \jugador\ deberá tocar para conseguir la puntuación deseada.
\item [Actores] \jugador.
\item [Precondiciones] Se ha elegido una canción.
\item [Postcondiciones] Se completa la interpretación de la canción, obteniendo una calificación
\item [Escenario principal] $\quad$
  \begin{enumerate}
  \item El \sistema\ carga la canción, leyendo las notas, y muestra en pantalla,
    mediante animaciones, el marcador de puntos y el pentagrama.
  \item El \sistema\ comienza a mostrar notas en el pentagrama, que van
    deslizándose hacia el lado izquierdo, en el que se encuentra la aguja de
    reproducción, e inicia el análisis del sonido.
  \item El \jugador\, al llegar la nota a la aguja de reproducción, toca la
    flauta con la altura y la duración correcta, de forma que el micrófono sea
    capaz de captar el sonido.
  \item El \sistema\ analiza el sonido que captura el micrófono y detecta la nota que toca el usuario.
  \item El \sistema\ determina que la nota es la correcta y suma los puntos correspondientes.
  \item Mientras existan más notas, se vuelve al punto 2.
  \item El \sistema\ determina que no hay más notas que mostrar, e inicia las
    animaciones para ocultar los elementos en pantalla.
  \item El \sistema\ pasa a la sección de \textit{Resultados de la interpretación}.
  \end{enumerate}
\item[Extensiones --- flujo alternativo] $\quad$
  \begin{description}

  \item [*a] El \jugador\ cierra la ventana.
    \begin{enumerate}
    \item El \sistema\ libera los recursos y sale de la aplicación.
    \end{enumerate}

  \item[*b] El \jugador\ pulsa la tecla \texttt{escape}.
    \begin{enumerate}
    \item El \sistema\ vuelve a la pantalla de \textit{selección de canción}.
    \end{enumerate}

  \item [3a] El \jugador\ toca el instrumento con intensidad insuficiente o nula
    y el sonido no llega al sistema.
    \begin{enumerate}
    \item El \sistema\ representa esta inconsistencia como un silencio.
    \end{enumerate}

  \item [4a] El \sistema\ no es capaz de determinar fehacientemente la nota que
    toca el usuario.
    \begin{enumerate}
    \item El \sistema\ representa esta inconsistencia como un silencio.
    \end{enumerate}

  \item[5a] El \sistema\ determina que la nota tocada por el usuario no es la
    que corresponde a la partitura.
    \begin{enumerate}
    \item El \sistema\ ignora esta situación y no suma los puntos al marcador.
    \end{enumerate}

  \end{description}  
\end{description}

\subsubsection{Caso de uso: resultados de la interpretación}

\begin{description}
\item [Descripción] Después de interpretar las notas de la partitura, se
  muestran los datos obtenidos del análisis de las notas tocadas por el \jugador.
\item [Actores] \jugador.
\item [Precondiciones] Se ha elegido e interpretado una canción.
\item [Postcondiciones] Se completa la partida actual.
\item [Escenario principal] $\quad$
  \begin{enumerate}
  \item El \sistema\ compara la puntuación conseguida con la máxima puntuación
    obtenible, y genera un porcentaje de aciertos.
  \item El \sistema\ muestra, mediante animaciones, un mensaje con información
    sobre la canción y sobre la interpretación del \jugador\, representada
    mediante un porcentaje de aciertos.
  \item El \sistema\ muestra un mensaje variable en función del número de
    aciertos conseguido.
  \item El \jugador\ revisa su puntuación y pulsa el botón \textit{volver} para
    ir de vuelta al menú de \textit{selección de canción}.
  \end{enumerate}
\item[Extensiones --- flujo alternativo] $\quad$
  \begin{description}

  \item [*a] El \jugador\ cierra la ventana.
    \begin{enumerate}
    \item El \sistema\ libera los recursos y sale de la aplicación.
    \end{enumerate}

  \item[4a] El \jugador\ pulsa la tecla \texttt{escape}.
    \begin{enumerate}
    \item El \sistema\ vuelve a la pantalla de \textit{selección de canción}.
    \end{enumerate}

  \end{description}  
\end{description}


\subsubsection{Caso de uso: analizador de notas}

\begin{description}
\item [Descripción] El \jugador\ elige la opción \textit{analizador de notas} en
  el menú principal y es llevado a esta sección, en la que el sistema
  representará gráficamente la nota que esté tocando con la flauta en cada
  instante, sin otra interacción
\item [Actores] \jugador.
\item [Precondiciones] Ninguna.
\item [Postcondiciones] Ninguna
\item [Escenario principal] $\quad$
  \begin{enumerate}
  \item El \jugador\ accede, desde el menú principal, al panel del analizador de notas.
  \item El \sistema\ muestra, mediante animaciones, la pantalla de la sección,
    representada mediante una fracción de partitura en la que se representará la
    nota que esté tocando el \jugador\ en cada momento.
  \item El \sistema\ inicia el análisis del sonido.
  \item El \jugador\ toca la nota que desee con su flauta, de forma que el
    micrófono sea capaz de captar el sonido.
  \item El \sistema\ analiza el sonido que captura el micrófono y detecta la
    nota que toca el usuario.
  \item El \sistema\ muestra en pantalla la nota, sobre la partitura,
    correspondiente a lo que ha tocado el usuario.
  \item Se repite el flujo desde el punto 4, mientras el \jugador\ no pulse en
    el botón volver.
  \item El \jugador\ pulsa en el botón \textit{volver}.
  \item El \sistema\ inicia las animaciones para ocultar los elementos en pantalla.
  \item El \sistema\ vuelve al menú principal.
  \end{enumerate}
\item[Extensiones --- flujo alternativo] $\quad$
  \begin{description}

  \item [*a] El \jugador\ cierra la ventana.
    \begin{enumerate}
    \item El \sistema\ libera los recursos y sale de la aplicación.
    \end{enumerate}

  \item[*b] El \jugador\ pulsa la tecla \texttt{escape}.
    \begin{enumerate}
    \item El \sistema\ vuelve al menú principal.
    \end{enumerate}

  \item [4a] El \jugador\ toca el instrumento con intensidad insuficiente o nula
    y el sonido no llega al sistema.
    \begin{enumerate}
    \item El \sistema\ representa esta inconsistencia como un silencio.
    \end{enumerate}

  \item [5a] El \sistema\ no es capaz de determinar fehacientemente la nota que
    toca el usuario.
    \begin{enumerate}
    \item El \sistema\ representa esta inconsistencia como un silencio.
    \end{enumerate}
  \end{description}  
\end{description}

\subsubsection{Caso de uso: calibración de micrófono}
\begin{description}
\item [Descripción] El \jugador\ elige la opción \textit{calibrar micrófono} en
  el menú principal y es llevado a esta sección, en la que el \sistema\ calibrará
  el micrófono de forma que sea posible aislar el sonido de la flauta del ruido
  ambiental.
\item [Actores] \jugador.
\item [Precondiciones] Ninguna.
\item [Postcondiciones] El \sistema\ obtiene un valor umbral con el que
  discernir entre el sonido del instrumento y el ruido ambiente.
\item [Escenario principal] $\quad$
  \begin{enumerate}
  \item El \jugador\ accede, desde el menú principal, al panel de calibración del micrófono.
  \item El \sistema\ muestra la sección, indicando con un mensaje que el usuario
    debe pulsar la tecla \texttt{escape} para iniciar la calibración.
  \item El \jugador\ pulsa la tecla \texttt{escape} y se mantiene en silencio.
  \item El \sistema\ inicia el análisis del sonido, guardando durante dos
    segundos los valores de ruido que lee del micrófono.
  \item El \sistema\ calcula, a partir de los valores leídos, el umbral de
    ruido, y muestra un mensaje informando del final del proceso.
  \item El \jugador\ pulsa la tecla \texttt{escape} y el \sistema\ vuelve al
    menú principal.
  \end{enumerate}
\item[Extensiones --- flujo alternativo] $\quad$
  \begin{description}

  \item [*a] El \jugador\ cierra la ventana.
    \begin{enumerate}
    \item El \sistema\ libera los recursos y sale de la aplicación.
    \end{enumerate}

  \item[*b] El \jugador\ pulsa la tecla \texttt{escape}.
    \begin{enumerate}
    \item El \sistema\ cancela la calibración y vuelve al menú principal.
    \end{enumerate}

  \item [5a] El \sistema\ encuentra valores inválidos al leer el ruido
    ambiental.
    \begin{enumerate}
    \item El \sistema\ informa al usuario del fallo del proceso de calibración.
    \item El \jugador\ pulsa la tecla \texttt{escape} y el \sistema\ vuelve al
      menú principal.
    \end{enumerate}
  \end{description}  
\end{description}

\subsubsection{Caso de uso: selección de lecciones}
\begin{description}
\item [Descripción] El \jugador\ elige la opción \textit{lecciones} en el menú
  principal y es llevado a esta sección, en la que el \sistema\ mostrará una
  lista de lecciones cargadas, entre las que el usuario deberá elegir.
\item [Actores] \jugador.
\item [Precondiciones] Ninguna.
\item [Postcondiciones] Se ha elegido una lección
\item [Escenario principal] $\quad$
  \begin{enumerate}
  \item El \jugador\ accede, desde el menú principal, al panel de selección de lecciones.
  \item El \sistema\ carga la lista de secciones y muestra, mediante
    animaciones, el panel, preseleccionando por defecto la primera lección.
  \item El \jugador\ utiliza los botones de la sección para elegir una de las
    lecciones, y activarla pulsando \textit{comenzar lección}.
  \item El \sistema\ oculta de forma animada el panel de selección de lecciones.
  \item El \sistema\ lee el fichero \texttt{xml} asociado a la lección indicada,
    cargando los elementos que la componen y las animaciones que se ejecutarán.
  \item El \sistema\ ejecuta las animaciones correspondientes a los elementos
    multimedia de la lección.
  \end{enumerate}
\item[Extensiones --- flujo alternativo] $\quad$
  \begin{description}

  \item [*a] El \jugador\ cierra la ventana.
    \begin{enumerate}
    \item El \sistema\ libera los recursos y sale de la aplicación.
    \end{enumerate}

  \item[2a] El \sistema\ detecta que una de las lecciones leídas no está
    correctamente construída.
    \begin{enumerate}
    \item El \sistema\ informa del error en el \textit{log} del programa y omite
      la carga de esa lección.
    \end{enumerate}

  \item[3a] El \jugador\ pulsa la tecla \texttt{escape}.
    \begin{enumerate}
    \item El \sistema\ vuelve al menú principal.
    \end{enumerate}

  \item[6a] El \jugador\ pulsa la tecla \texttt{escape}.
    \begin{enumerate}
    \item El \sistema\ vuelve al menú de selección de lecciones.
    \end{enumerate}

  \end{description}  
\end{description}

\section{Modelo conceptual de datos}

El modelo conceptual de datos representa, de forma esquemática, las clases que
modelan el sistema y las relaciones que existen entre ellas, además de una
pequeña introducción a su utilidad. 

\begin{description}
\item[Juego] Clase de control general. Gestiona el flujo de ejecución principal,
  así como de la gestión de estados, que permite pasar de una sección a otra del
  juego.
\item[Estado] Clase base para los diferentes estados del juego. Las clases
  correspondientes a las secciones se basarán en esta clase para interactuar con
  el gestor de estados y poder pasar de una parte del juego a otra.
\item[EstadoMenú] Representa el estado de juego para el menú principal, desde el
  que se accede al resto de opciones del juego. 
\item[EstadoAnalizador] Representa el estado del analizador básico de
  notas. Contendrá los elementos necesarios para iniciar el análisis del audio,
  así como los elementos de la interfaz.
\item[EstadoCalibrarMicro] Representa el estado en el que se calibra el
  micrófono. Al igual que la clase \textit{EstadoAnalizador}, deberá ser capaz
  de acceder al sistema de audio para poder leer el volumen ambiente y así
  calibrar el micrófono.
\item[EstadoImagenFija] Modela una imagen fija a modo de pantalla de créditos,
  de forma que sea sencillo añadir imágenes al inicio del juego, como firmas de
  desarrolladores, logotipos de patrocinadores, etcétera.
\item[EstadoMenúCanciones] Comprende el menú de selección de canciones, que se
  encargará de leer los ficheros de canciones disponibles. Además, también se
  encargará de lanzar las canciones en forma de estados secundarios.
\item[EstadoCanción] Corresponde a la canción que se va a interpretar, lanzada desde
  el estado \textit{EstadoMenuCanciones}.
\item[EstadoMenúLecciones] Corresponde al menú de elección de lecciones, que
  leerá y listará los ficheros de lección disponibles, y se encargará de lanzar
  la lección elegida.
\item[EstadoLección] Corresponde a la canción elegida desde el menú
  \textit{EstadoMenúLecciones}.
\item[Analizador] Controla la gestión del subsistema de audio y el análisis de
  la entrada. Deberá proporcionar información sobre el volumen de la entrada
  (para la calibración del micrófono) así como de la nota detectada en cada
  instante.
\item[Animación] Se encargará de facilitar la creación de animaciones en forma
  de interpolación de valores, válidas para cambios de posición, opacidad,
  etcétera.
\item[Elemento] Esta clase de ayuda facilitará la carga y dibujado de elementos
  para la interfaz, además de servir de capa de abstracción para las
  animaciones.
\item[ElementoImagen] Especialización de la clase \textit{Elemento} para
  imágenes.
\item[ElementoTexto] Especialización de la clase \textit{Elemento} para textos.
\item[ElementoCombinado] Especialización de la clase \textit{Elemento} que
  combina imagen y texto, a usar en casos como los botones del menú.
\item[SistemaPartículas] Representa un sistema de partículas simple, para
  generar efectos de destellos y fuegos artificiales.
\item[Nota] Simboliza cada una de las notas cargadas que componen una canción.
\end{description}

\begin{figure}[htp!]
  \centering
  \includegraphics[angle=90]{4_analisis/imagen_diagrama_clases_conceptuales}
  \caption{Diagrama de clases conceptuales}
\end{figure}

\pagebreak


\section{Modelo de comportamiento del sistema}
En esta sección vamos a especificar cómo se comporta el sistema en forma de dos
elementos fundamentales.
\begin{itemize}
\item En primer lugar, los \textbf{diagramas de secuencia} mostrarán el flujo de
  eventos entre los actores que participan en la aplicación.
\item En segundo lugar, los \textbf{contratos de las operaciones} detallarán las
  condiciones y efectos que tendrán lugar al ejecutarse las operaciones en el
  sistema.
\end{itemize}

\begin{nota}
  No se han reflejado, por triviales, los escenarios alternativos en los que el
  usuario cierra la ventana, correspondientes a los flujos \textit{*a} definidos
  en la sección anterior.
\end{nota}

\subsection{Caso de uso: inicio del juego}

\subsubsection{Escenario principal}

\begin{figure}[h!]
  \centering
  \includegraphics[trim=0cm 12cm 0cm 0cm, clip=true, width=0.5\textwidth]{4_analisis/diagsec_caso1_esc1}
  \caption{Diagrama de secuencia, incio del juego, escenario principal}
\end{figure}

\begin{description}
\item[Operación] IniciarJuego()
\item[Actores] \jugador\, \sistema\
\item[Responsabilidades] Cargar y lanzar la aplicación, mostrar los títulos de
  crédito y el menú principal.
\item[Precondiciones] Ninguna.
\item[Postcondiciones] $\quad$

  \begin{itemize}
  \item Se crea una instancia de la clase \textit{Juego}, que gestiona la
    creación y destrucción de los estados.
  \item Se crean y posteriormente destruyen dos estados
    \textit{EstadoImagenFija} para mostrar los títulos de crédito.
  \item Se crea y permanece un estado \textit{EstadoMenú}, que representa el
    menú principal de la aplicación.
  \end{itemize}

\end{description}

\begin{description}
\item[Operación] SeleccionarSeccionCanciones()
\item[Actores] \jugador\, \sistema\
\item[Responsabilidades] Esconder el menú principal y cargar el menú de
  selección de canciones.
\item[Precondiciones] $\quad$

  \begin{itemize}
  \item El estado actual es una instancia de \textit{EstadoMenú}.
  \end{itemize}

\item[Postcondiciones] Se destruye el estado \textit{EstadoMenú} y se carga
  \textit{EstadoMenúCanciones}.
\end{description}

\subsubsection{Escenario alternativo 4a}
\begin{figure}[h!]
  \centering
  \includegraphics[trim=0cm 12cm 0cm 0cm, clip=true, width=0.5\textwidth]{4_analisis/diagsec_caso1_esc2}
  \caption{Diagrama de secuencia, incio del juego, escenario alternativo 4a}
\end{figure}

\begin{description}
\item[Operación] SeleccionarSeccionAnalizador()
\item[Actores] \jugador\, \sistema\
\item[Responsabilidades] Esconder el menú principal y cargar la sección de
  análisis de notas.
\item[Precondiciones] $\quad$
  \begin{itemize}
  \item El estado actual es una instancia de \textit{EstadoMenú}.
  \end{itemize}
\item[Postcondiciones] Se destruye el estado \textit{EstadoMenú} y se carga
  \textit{EstadoAnalizador}.
\end{description}

\subsubsection{Escenario alternativo 4b}
\begin{figure}[h!]
  \centering
  \includegraphics[trim=0cm 12cm 0cm 0cm, clip=true, width=0.5\textwidth]{4_analisis/diagsec_caso1_esc3}
  \caption{Diagrama de secuencia, incio del juego, escenario alternativo 4b}
\end{figure}

\begin{description}
\item[Operación] SeleccionarSeccionLecciones()
\item[Actores] \jugador\, \sistema\
\item[Responsabilidades] Esconder el menú principal y cargar la menú de
  selección de lecciones.
\item[Precondiciones] $\quad$
  \begin{itemize}
  \item El estado actual es una instancia de \textit{EstadoMenú}.
  \end{itemize}
\item[Postcondiciones] Se destruye el estado \textit{EstadoMenú} y se carga
  \textit{EstadoMenúLecciones}.
\end{description}

\subsubsection{Escenario alternativo 4c}
\begin{figure}[h!]
  \centering
  \includegraphics[trim=0cm 12cm 0cm 0cm, clip=true, width=0.5\textwidth]{4_analisis/diagsec_caso1_esc4}
  \caption{Diagrama de secuencia, incio del juego, escenario alternativo 4c}
\end{figure}

\begin{description}
\item[Operación] SeleccionarSeccionCalibracion()
\item[Actores] \jugador\, \sistema\
\item[Responsabilidades] Esconder el menú principal y cargar la sección de
  calibración de micrófono.
\item[Precondiciones] $\quad$
  \begin{itemize}
  \item El estado actual es una instancia de \textit{EstadoMenú}.
  \end{itemize}
\item[Postcondiciones] Se destruye el estado \textit{EstadoMenú} y se carga
  \textit{EstadoCalibrarMicro}.
\end{description}

\subsubsection{Escenario alternativo 4d}
\begin{figure}[h!]
  \centering
  \includegraphics[trim=0cm 12cm 0cm 0cm, clip=true, width=0.5\textwidth]{4_analisis/diagsec_caso1_esc5}
  \caption{Diagrama de secuencia, incio del juego, escenario alternativo 4d}
\end{figure}

\begin{description}
\item[Operación] SeleccionarSalir()
\item[Actores] \jugador\, \sistema\
\item[Responsabilidades] Esconder el menú principal, descargar los recursos y
  cerrar la aplicación.
\item[Precondiciones] $\quad$
  \begin{itemize}
  \item El estado actual es una instancia de \textit{EstadoMenú}.
  \end{itemize}
\item[Postcondiciones] Se destruye el estado \textit{EstadoMenú}, se destruye la
  instancia de la clase \textit{Juego} y termina la ejecución de la aplicación.
\end{description}

\subsection{Selección de canción}

\subsubsection{Escenario principal}
\begin{figure}[h!]
  \centering
  \includegraphics[trim=0cm 12cm 0cm 0cm, clip=true, width=0.5\textwidth]{4_analisis/diagsec_caso2_esc1}
  \caption{Diagrama de secuencia, selección de canción, escenario principal}
\end{figure}

\begin{description}
\item[Operación] ListarCanciones()
\item[Actores] \jugador\, \sistema\
\item[Responsabilidades] Cargar y mostrar la lista de canciones cargadas en el
  sistema.
\item[Precondiciones] Se ordenó la carga del estado \textit{EstadoMenuCanción}
\item[Postcondiciones] $\quad$
  \begin{itemize}
  \item El estado actual es una instancia de \textit{EstadoMenuCanción}.
  \item Se ha cargado la lista de canciones y se muestra en pantalla.
  \end{itemize}
\end{description}

\begin{description}
\item[Operación] ElegirCanción()
\item[Actores] \jugador\, \sistema\
\item[Responsabilidades] Cargar la canción que el usuario ha elegido para
  interpretar.
\item[Precondiciones] Existe una lista de canciones cargada de entre las que el
  usuario ha elegido una.
\item[Postcondiciones] $\quad$
  \begin{itemize}
  \item Se carga la canción indicada.
  \item Se oculta la lista de canciones.
  \item Se pasa a un sub-estado de interpretación de canción.
  \end{itemize}
\end{description}
\subsubsection{Escenario alternativo 3a}
\begin{figure}[h!]
  \centering
  \includegraphics[trim=0cm 12cm 0cm 0cm, clip=true, width=0.5\textwidth]{4_analisis/diagsec_caso2_esc2}
  \caption{Diagrama de secuencia, selección de canción, escenario alternativo 3a}
\end{figure}
\begin{description}
\item[Operación] ElegirVolver()
\item[Actores] \jugador\, \sistema\
\item[Responsabilidades] Descargar la sección actual y volver al menú anterior.
\item[Precondiciones] El estado actual es una instancia de \textit{EstadoMenuCanción}.
\item[Postcondiciones] $\quad$
  \begin{itemize}
  \item El estado instancia de \textit{EstadoMenuCanción} queda descargado.
  \item Se carga y se muestra \textit{EstadoMenú}.
  \end{itemize}
\end{description}

\subsection{Interpretación de canción}
\begin{nota}
  No se reflejan los escenarios alternativos al estar englobados en la operación
  \textit{InteractuarConFlauta}.
\end{nota}

\subsubsection{Escenario principal}
\begin{figure}[h!]
  \centering
  \includegraphics[trim=0cm 8cm 0cm 0cm, clip=true, width=0.5\textwidth]{4_analisis/diagsec_caso3}
  \caption{Diagrama de secuencia, interpretación de canción, escenario principal}
\end{figure}

\begin{description}
\item[Operación] IniciarInterpretación()
\item[Actores] \jugador\, \sistema\
\item[Responsabilidades] Parsear el fichero de canción, cargar la interfaz y
  comenzar la interpretación.
\item[Precondiciones] El usuario ha elegido una canción en el estado anterior.
\item[Postcondiciones] $\quad$
  \begin{itemize}
  \item Se muestra la interfaz de interpretación de canción.
  \item El fichero de canción queda cargado e interpretado, instanciando los
    elementos de la clase \textit{Nota} que sean necesarios.
  \item Comienza la interpretación
  \end{itemize}
\end{description}

\begin{description}
\item[Operación] InteractuarConFlauta()
\item[Actores] \jugador\, \sistema\
\item[Responsabilidades] El \jugador\ interactúa con el sistema mediante la
  flauta a través del micrófono, y el \sistema\ analiza los datos y muestra una
  respuesta en pantalla.
\item[Precondiciones] $\quad$
  \begin{itemize}
  \item La interpretación ha comenzado.
  \item El micrófono está correctamente configurado.
  \end{itemize}
\item[Postcondiciones] $\quad$
  \begin{itemize}
  \item El sistema captura y analiza los datos de audio.
  \item Según el análisis, el sistema responde de una forma u otra (según los
    escenarios alternativos 3a, 4a y 5a del caso de uso \textit{interpretación
      de canción}).
  \end{itemize}
\end{description}

\begin{description}
\item[Operación] FinalizarInterpretación()
\item[Actores] \jugador\, \sistema\
\item[Responsabilidades] Descargar la pantalla de interpretación, descargar la
  canción y finalizar la interpretación.
\item[Precondiciones] Todas las notas se han interpretado.
\item[Postcondiciones] $\quad$
  \begin{itemize}
  \item Se descarga la \textit{Canción} actual.
  \item Se ocultan los elementos de la interfaz de interpretación.
  \item Se lanza la sección de puntuación.
  \end{itemize}
\end{description}


\subsection{Resultados de interpretación}

\subsubsection{Escenario principal}
\begin{figure}[h!]
  \centering
  \includegraphics[trim=0cm 12cm 0cm 0cm, clip=true, width=0.5\textwidth]{4_analisis/diagsec_caso4}
  \caption{Diagrama de secuencia, resultados de interpretación, escenario principal}
\end{figure}

\begin{description}
\item[Operación] MostrarResultados()
\item[Actores] \jugador\, \sistema\
\item[Responsabilidades] Interpretar los resultados de la interpretación y
  mostrar los resultados en pantalla.
\item[Precondiciones] El \jugador\ ha concluído satisfactoriamente una
  interpretación completa de una canción, obteniendo una suma de puntos $X$.
\item[Postcondiciones] Mostrar en pantalla los resultados en forma de porcentaje
  de aciertos, y un mensaje según aquél.
\end{description}

\begin{description}
\item[Operación] ElegirVolver()
\item[Actores] \jugador\, \sistema\
\item[Responsabilidades] Descargar la sección actual y volver al menú anterior.
\item[Precondiciones] $\quad$
  \begin{itemize}
  \item La aplicación se encuentra en la pantalla de muestra de resultados.
  \item El usuario ha pulsado la tecla \texttt{escape} o el botón \textit{volver}.
  \end{itemize}
\item[Postcondiciones] $\quad$
  \begin{itemize}
  \item Se descargan todos los datos referentes a la canción actual.
  \item Se carga y se muestra \textit{EstadoMenúCanciones}.
  \end{itemize}
\end{description}

\subsection{Analizador de notas}

\begin{nota}
  No se reflejan los escenarios alternativos al estar englobados en la operación
  \textit{InteractuarConFlauta}.
\end{nota}

\subsubsection{Escenario principal}
\begin{figure}[h!]
  \centering
  \includegraphics[trim=0cm 8cm 0cm 0cm, clip=true, width=0.5\textwidth]{4_analisis/diagsec_caso5}
  \caption{Diagrama de secuencia, interpretación de canción, escenario principal}
\end{figure}

\begin{description}
\item[Operación] IniciarAnálisis
\item[Actores] \jugador\, \sistema\
\item[Responsabilidades] Cargar la interfaz e iniciar el análisis de notas.
\item[Precondiciones] El usuario eligió la sección \textit{Analizador de notas}
  en el menú principal.
\item[Postcondiciones] $\quad$
  \begin{itemize}
  \item Aparece la interfaz del analizador de notas.
  \item Se inicia el análisis de notas
  \end{itemize}
\end{description}

\begin{description}
\item[Operación] InteractuarConFlauta
\item[Actores] \jugador\, \sistema\
\item[Responsabilidades] El \jugador\ toca notas en la flauta y el \sistema\
  captura y reconoce el audio, indicando la nota tocada en pantalla.

\item[Precondiciones] Se ha iniciado el análisis.

\item[Postcondiciones] $\quad$
  \begin{itemize}
  \item El \sistema\ recoge y analiza el sonido que emite la flauta del \jugador.
  \item El \sistema\ representa en pantalla la nota identificada, o no muestra
    nada en caso de identificación defectuosa.
  \end{itemize}
\end{description}

% \begin{description}
% \item[Operación] ElegirVolver()
% \item[Actores] \jugador\, \sistema\
% \item[Responsabilidades] Descargar la sección actual y volver al menú anterior.
% \item[Precondiciones] $\quad$
%   \begin{itemize}
%   \item El estado actual es una instancia de \textit{EstadoAnalizador}.
%   \item Hay un análisis en curso.  
%   \end{itemize}
  
% \item[Postcondiciones] $\quad$
%   \begin{itemize}
%   \item Concluye el análisis actual.
%   \item Se destruye el estado.
%   \item Se carga y se muestra \textit{EstadoMenú}.
%   \end{itemize}
% \end{description}

\subsection{Calibración de micrófono}

\subsubsection{Escenario principal}
\begin{figure}[h!]
  \centering
  \includegraphics[trim=0cm 8cm 0cm 0cm, clip=true, width=0.5\textwidth]{4_analisis/diagsec_caso6_esc1}
  \caption{Diagrama de secuencia, calibración de micrófono, escenario principal}
\end{figure}

\begin{description}
\item[Operación] CargarCalibración()
\item[Actores] \jugador\, \sistema\
\item[Responsabilidades] Cargar la sección y preparar el sistema para comenzar
  la calibración del micrófono.
\item[Precondiciones] El usuario ha elegido en el menú principal la opción
  \textit{Calibrar micrófono}.
\item[Postcondiciones] La sección está cargada y la calibración lista para
  iniciarse.
\end{description}

\begin{description}
\item[Operación] CalibrarMicrófono()
\item[Actores] \jugador\, \sistema\
\item[Responsabilidades] Llevar a cabo la calibración correcta del micrófono.
\item[Precondiciones] El usuario ha lanzado la calibración del micrófono.
\item[Postcondiciones] $\quad$
  \begin{itemize}
  \item Se cierra el sistema de sonido.
  \item Se obtiene un valor umbral de ruido ambiente, fruto de una calibración
    exitosa.
  \end{itemize}
\end{description}

\begin{description}
\item[Operación] ElegirVolver()
\item[Actores] \jugador\, \sistema\
\item[Responsabilidades] Descargar la sección actual y volver al menú anterior.
\item[Precondiciones] $\quad$
  \begin{itemize}
  \item La calibración ha concluído exitosamente.
  \item El usuario ha pulsado la tecla \texttt{escape} o el botón \textit{volver}.
  \end{itemize}
\item[Postcondiciones] $\quad$
  \begin{itemize}
  \item Se descarga la sección actual.
  \item Se carga y se muestra \textit{EstadoMenúCanciones}.
  \end{itemize}
\end{description}

\subsubsection{Escenario alternativo 5a}
\begin{figure}[h!]
  \centering
  \includegraphics[trim=0cm 8cm 0cm 0cm, clip=true, width=0.5\textwidth]{4_analisis/diagsec_caso6_esc2}
  \caption{Diagrama de secuencia, calibración de micrófono, escenario principal}
\end{figure}
\begin{description}
\item[Operación] CalibrarMicrófono()
\item[Actores] \jugador\, \sistema\
\item[Responsabilidades] Llevar a cabo la calibración correcta del micrófono.
\item[Precondiciones] El usuario ha lanzado la calibración del micrófono.
\item[Postcondiciones] $\quad$
  \begin{itemize}
  \item Se cierra el sistema de sonido.
  \item La calibración ha fallado.
  \item No se obtiene valor umbral de ruido ambiental.
  \end{itemize}
\end{description}

\subsection{Selección de lecciones}

\subsubsection{Escenario principal}
\begin{figure}[h!]
  \centering
  \includegraphics[trim=0cm 4cm 0cm 0cm, clip=true, width=0.5\textwidth]{4_analisis/diagsec_caso7_esc1}
  \caption{Diagrama de secuencia, selección delecciones, escenario principal}
\end{figure}


\begin{description}
\item[Operación] ListarLecciones()
\item[Actores] \jugador\, \sistema\
\item[Responsabilidades] Mostrar el menú de selección de lecciones y listar
  todas las lecciones disponibles.
\item[Precondiciones] El usuario eligió la opción \textit{Lecciones} en el menú
  principal.
\item[Postcondiciones] $\quad$
  \begin{itemize}
  \item Se muestra la interfaz del menú de selección de lecciones.
  \item Se listan las lecciones cargadas en el sistema
  \end{itemize}
\end{description}

\begin{description}
\item[Operación] ElegirLección()
\item[Actores] \jugador\, \sistema\
\item[Responsabilidades] Cargar y mostrar la lección elegida por el usuario.
\item[Precondiciones] El menú de selección de lecciones está cargado y el
  usuario ha elegido una de las lecciones.
\item[Postcondiciones]$\quad$
  \begin{itemize}
  \item Se oculta el menú de selección de lecciones.
  \item Se interpreta el fichero de lección elegido.
  \item Se muestran los elementos multimedia pertenecientes a la lección
    elegida.
  \end{itemize}
\end{description}

\begin{description}
\item[Operación] VolverMenuLecciones()
\item[Actores] \jugador\, \sistema\
\item[Responsabilidades] Descargar la lección actual y volver al menú anterior.
\item[Precondiciones] $\quad$
  \begin{itemize}
  \item El usuario ha terminado de ver la lección elegida.
  \item El usuario ha pulsado la tecla \texttt{escape} o el botón \textit{volver}.
  \end{itemize}
\item[Postcondiciones] $\quad$
  \begin{itemize}
  \item Se descarga la lección actual.
  \item Se muestra de nuevo el menú de selección de lecciones.
  \end{itemize}
\end{description}

\begin{description}
\item[Operación] VolverMenuPrincipal()
\item[Actores] \jugador\, \sistema\
\item[Responsabilidades] Descargar el menú de selección de lecciones y volver al
  menú principal.
\item[Precondiciones] El usuario ha pulsado la tecla \texttt{escape} o el botón \textit{volver}.
\item[Postcondiciones] $\quad$
  \begin{itemize}
  \item Se descarga el menú de selección de lecciones.
  \item Se carga y muestra el menú principal
  \end{itemize}
\end{description}

\subsubsection{Escenario alternativo}
\begin{figure}[h!]
  \centering
  \includegraphics[trim=0cm 12cm 0cm 0cm, clip=true, width=0.5\textwidth]{4_analisis/diagsec_caso7_esc2}
  \caption{Diagrama de secuencia, selección de lecciones, escenario alternativo}
\end{figure}

\begin{description}
\item[Operación] ListarLecciones()
\item[Actores] \jugador\, \sistema\
\item[Responsabilidades] Mostrar el menú de selección de lecciones y listar
  todas las lecciones disponibles.
\item[Precondiciones] El usuario eligió la opción \textit{Lecciones} en el menú
  principal.
\item[Postcondiciones] $\quad$
  \begin{itemize}
  \item Se muestra la interfaz del menú de selección de lecciones.
  \item Se listan las lecciones cargadas en el sistema
  \end{itemize}
\end{description}

\begin{description}
\item[Operación] VolverMenuPrincipal()
\item[Actores] \jugador\, \sistema\
\item[Responsabilidades] Descargar el menú de selección de lecciones y volver al
  menú principal.
\item[Precondiciones] El usuario ha pulsado la tecla \texttt{escape} o el botón \textit{volver}.
\item[Postcondiciones] $\quad$
  \begin{itemize}
  \item Se descarga el menú de selección de lecciones.
  \item Se carga y muestra el menú principal
  \end{itemize}
\end{description}

%%% Local Variables: 
%%% mode: latex
%%% TeX-master: "../memoria"
%%% End: 

%%% Local Variables: 
%%% mode: latex
%%% TeX-master: "../memoria"
%%% End: 


\chapter{Diseño}
\label{chap:diseno}
%En este capítulo se presentarán los detalles de diseño del proyecto, basándonos
en el análisis mostrado en el capítulo anterior. El modelo de clases de diseño
representado aquí es más fiel a la implementación final que el diagrama de
clases conceptuales, pero aún así hay detalles que no se han contemplado por ser
demasiado cercanos a los detalles de implementación.

También en el presente capítulo se detallarán las decisiones de diseño en
relación al aspecto visual de la aplicación, extendiéndonos en el proceso de
creación del logotipo del juego así como de la interfaz gráfica de usuario.

\section{Diagrama de clases de diseño}

Tras la fase de diseño, componían el sistema más de 40 clases, por lo que hemos
tenido que dividir los diagramas en varias partes, intentando seguir cierto
criterio a la hora de elegir qué clases formarán parte de cada división.

En todos los diagramas aparecen, en aras de mantener el contexto, las clases
básicas de la aplicación: \textit{Juego} y \textit{Estado}. Además, hay algunas
otras clases que también se repetirán entre diagramas por conveniencia. Para
mantener la legibilidad, se han ocultado los miembros privados y protegidos.

\begin{itemize}
\item En el primer diagrama (figura~\ref{fig:diagrama_clases_1}) aparecen las clases
  relacionadas con el \textit{menú principal}, clases de utilidades (logging y
  animación), y clases para representar elementos gráficos.
\item En el segundo diagrama (figura~\ref{fig:diagrama_clases_2}) aparecen las
  clases relacionadas con el subsistema de análisis del audio, así como las
  secciones \textit{Analizador de Notas} y \textit{Calibrar micrófono}.
\item En el tercer diagrama (figura~\ref{fig:diagrama_clases_3}) aparecen todas las
  clases relacionadas con la sección de \textit{Canciones}.
\item En el cuarto y último diagrama (figura~\ref{fig:diagrama_clases_4}) aparecen
  las clases relacionadas con el motor de \textit{lecciones}.
\end{itemize}

\begin{figure}[htp!]
  \centering
  \includegraphics[width=\textwidth]{5_diseno/diagrama1}
  \caption{Diagrama de clases de diseño, parte I}
  \label{fig:diagrama_clases_1}
\end{figure}

\begin{figure}[htp!]
  \centering
  \includegraphics[width=\textwidth, clip=true, trim=0cm 0cm 0cm 0cm]{5_diseno/diagrama2}
  \caption{Diagrama de clases de diseño, parte II}
  \label{fig:diagrama_clases_2}
\end{figure}

\begin{figure}[htp!]
  \centering
  \includegraphics[width=\textwidth, clip=true, trim=0cm 0cm 0cm 0cm]{5_diseno/diagrama3}
  \caption{Diagrama de clases de diseño, parte III}
  \label{fig:diagrama_clases_3}
\end{figure}

\begin{figure}[htp!]
  \centering
  \includegraphics[width=\textwidth, clip=true, trim=0cm 0cm 0cm 0cm]{5_diseno/diagrama4}
  \caption{Diagrama de clases de diseño, parte IV}
  \label{fig:diagrama_clases_4}
\end{figure}

\section{Diseño visual}
Teniendo muy presente que el público objetivo de \textbf{oFlute} es joven y
dinámico, el diseño visual del proyecto intentó adaptarse a lo que la audiencia
podría considerar más atractivo. Desde el inicio se fijaron ciertas premisas o
\textit{guías de branding} para oFlute, que se siguieron a rajatabla a la hora
de diseñar cada uno de los elementos.

\textbf{Interfaces limpias y minimalistas}. Era imprescindible que las
interfaces gráficas de cada una de las secciones tuviera un diseño limpio y
cuidado, manteniendo en el mínimo la cantidad de elementos a mostrar en pantalla
sin descuidar el diseño.

Para conseguirlo, se optó por utilizar un fondo blanco con un sutil patrón gris
claro, que se mantiene a lo largo de todas las secciones. Un gran ejemplo de esto
es la imagen de los títulos de crédito del juego.

\begin{figure}[htp!]
  \centering
  \includegraphics[width=0.7\textwidth]{apendice_manual_usuario/imagen_estadoIntro}
  \caption{Imagen de títulos de crédito de oFlute}
\end{figure}

\textbf{Paleta de colores pastel}. Para acompañar al blanco limpio de los
fondos de la interfaz se ideó una paleta de colores amplia, basado en tonos
brillantes, cercanos a pastel. Esta paleta está presente tanto en el logotipo de
oFlute, como en el resto de secciones. En especial, el menú principal se
benefició enormemente del uso de esta paleta de colores en los botones de las
secciones.

\begin{figure}[htp!]
  \centering
  \includegraphics[width=0.8\textwidth]{5_diseno/imagen_escala_colores}
  \caption{Detalle de la paleta de colores de oFlute}
\end{figure}

\textbf{Logotipo}. A la hora de diseñar el logtipo, se hizo un
\textit{brainstorming} inicial en busca de conceptos relacionados con el
proyecto: notas musicales, flautas, pentagramas, claves de sol, etcétera.

Finalmente, decidimos utilizar el concepto de la flauta a través del espacio
negativo que genera dibujar solo los orificios de la misma, manteniendo el
tamaño original de los mismos.

\begin{figure}[htp!]
  \centering
  \includegraphics[width=\textwidth]{5_diseno/imagen_composite}
  \caption{Logotipo de oFlute original sobre blanco, sobre negro y en blanco y
    negro}
\end{figure}

\textbf{Tipografía}. Para adecuarse al diseño minimalista propuesto, se buscaron
tipografías simples y limpias. Para el logotipo se usó la fuente 232 MKSD
Round~\cite{fuentelogo}, y para los textos de la aplicación se usó la fuente
Miso~\cite{fuentetexto}.

%%% Local Variables: 
%%% mode: latex
%%% TeX-master: "../memoria"
%%% End: 


\chapter{Implementación}
\label{chap:implementacion}
%Como complemento a la lectura de este capítulo se recomienda tener una copia
local del repositorio del proyecto, disponible para su libre descarga desde la
forja oficial~\cite{forja}.

\section{Entorno de construcción}

En esta sección se detalla el marco tecnológico utilizado para al construcción
del sistema, basándonos en lo presentado en la sección~\ref{sec:alternativas-solucion},~\nameref{sec:alternativas-solucion}.

\subsection{Aplicación web}

El entorno tecnológico de la aplicación web ha resultado ser muy variado y rico
en tecnologías en total vigencia en la actualidad. 

\subsubsection{Herramientas de diseño y desarrollo}

Como editor principal se ha utilizado \textbf{Sublime Text
  2}~\cite{sublime-text}, un editor freeware escrito en Python especialmente
dotado para el desarrollo de aplicaciones web. Cuenta con un gran número de
características que lo han hecho destacar en los últimos años, como la habilidad
de seleccionar y editar a la vez varios fragmentos de texto, el rápido motor de
\textit{búsqueda difusa} para encontrar ficheros y cadenas o su capacidad de
extensión a través de paquetes escritos por terceros y fácilmente instalables.

Para ciertas tareas, como la escritura de la presente memoria, se ha hecho uso
del veterano editor \textbf{Emacs}~\cite{emacs}, que sigue siendo la mejor
opción a la hora de editar documentos de \LaTeX. 

Se ha utilizado \textbf{draw.io}~\cite{draw.io} para la generación de diagramas
de casos de uso, secuencia, diagramas de flujo y similares. Se trata de una
herramienta web, actualmente integrable en Google Drive, con un enorme número de
elementos de dibujo y la capacidad de exportar en un gran número de formatos,
entre ellos en formato vectorial \ac{PDF}. Para el diseño visual de las
interfaces de usuario se ha utilizado \textbf{Adobe Photoshop}~\cite{photoshop}.

\subsubsection{Gestión de dependencias}

Para facilitar la gestión de las dependencias del proyecto se han utilizado
\textbf{VirtualEnv}~\cite{virtualenv} y
\textbf{VirtualEnvWrapper}~\cite{virtualenvwrapper}. Estas herramientas permiten
generar \textit{entornos virtuales} para cada proyecto, en los que se instalan
las dependencias necesarias. Estos entornos se activan y desactivan, de forma
que las bibliotecas instaladas en el entorno virtual de un proyecto no son
accesibles desde el entorno del otro proyecto. Esto evita la polución del nivel
general de bibliotecas y facilita el control estricto de las versiones de las
dependencias. 

Junto a virtualenv, la herramienta \textbf{pip}~\cite{pip} permite guardar en un
fichero anexo la lista de dependencias de un proyecto, de forma que sea fácil
reinstalarlas todas si hubiese que repetir la instalación en otro sistema.

\subsubsection{Control de versiones}

Todo el código fuente del proyecto se encuentra alojado en un repositorio
público en \textbf{GitHub}~\cite{forja}, haciendo uso de los planes
gratuitos. GitHub es una forja que utiliza el sistema de control de versiones
\textbf{Git}~\cite{git}. Además del alojamiento de código, GitHub provee
numerosas funcionalidades adicionales, tanto a nivel social (permitiendo a las
forjas tener \textit{followers}, por ejemplo) como a nivel funcional (ofreciendo
sistemas de tickets, estadísticas, etcétera).

El uso de un control de versiones es fundamental por varios motivos. En primer
lugar, sirve como sistema de copia de seguridad. En segundo lugar, permite
deshacer cambios en el código que no funcionen bien, siendo siempre posible
\textit{volver atrás}. Por último, sirve como \textit{cuaderno de bitácora}
improvisado, ya que se guarda el historial de \textit{commits} que el
desarrollador va enviando junto a los mensajes, siendo posible ver en una línea
temporal el progreso del trabajo.

\subsubsection{Lenguaje de programación}

Como se ha comentado en numerosas ocasiones en la presente memoria, el lenguaje
de programación elegido para el desarrollo del proyecto es
\textbf{Python}~\cite{Python}, un lenguaje interpretado de alto nivel
desarrollado por Guido Van Rossum en 1991. Python soporta múltiples paradigmas
de programación, desde la orientación a objetos hasta la programación funcional,
pasando por el clásico estilo imperativo. Su principal uso ha sido como lenguaje
de scripting, pero también tiene su hueco en contextos más amplios como lenguaje
principal.

Su facilidad de aprendizaje, lo simple de su sintaxis (basada en la indentación
para marcar los bloques) y su extensibilidad (sobre todo gracias al uso de
métodos \textit{mágicos} y metaprogramación) han hecho que el lenguaje sea muy
popular y su uso en los últimos años se haya expandido enormemente.

\subsubsection{Framework de desarrollo}

El framework de desarrollo elegido es \textbf{Django}~\cite{Django}, un
proyecto de código abierto nacido en 2005. La meta fundamental de Django es
facilitar la creación de sitios web complejos, haciendo especial hincapié en
mantener un ritmo de desarrollo rápido y evitar la duplicidad de código. Así,
gran parte del potencial de Django es la gran cantidad de funcionalidad que ya
trae integrada, ya sea en el propio núcleo del framework o a través de
\textit{apps} oficiales que se incluyen en la distribución oficial.

Los elementos más destacables de Django son:

\begin{itemize}
\item Un sistema de mapeo objeto-relacional, que facilita enormemente el trabajo
  con elementos de bases de datos a través de instancias de modelos.
\item Un sistema de procesamiento de peticiones a través de vistas.
\item Un despachador de URLs basado en expresiones regulares.
\item Soporte para plantillas.
\item Un sistema de autenticación extensible.
\item Una interfaz de administración que se adapta dinámicamente a cualquier proyecto.
\item Un sistema de comentarios.
\item Soporte integrado contra numerosos tipos de ataques vía web, como
  inyección SQL o cross-site scripting.
\end{itemize}

\subsubsection{Nivel de persistencia}

\textbf{SQLite} es un sistema de gestión de bases de datos relacional de dominio público,
contenido en una pequeña biblioteca. A diferencia de los sistema de gestión de
bases de datos cliente-servidor, el motor de SQLite no es un proceso
independiente con el que el programa principal se comunica. En lugar de eso, la
biblioteca SQLite se enlaza con el programa pasando a ser parte integral del
mismo. El programa utiliza la funcionalidad de SQLite a través de llamadas
simples a subrutinas y funciones. Esto reduce la latencia en el acceso a la base
de datos, debido a que las llamadas a funciones son más eficientes que la
comunicación entre procesos. 

El conjunto de la base de datos (definiciones, tablas, índices, y los propios
datos), son guardados como un sólo fichero estándar en la máquina host. Este
diseño simple se logra bloqueando todo el fichero de base de datos al principio
de cada transacción. Esto no impide que varios procesos o hilos pueden acceder a
la misma base de datos sin problemas.

En general, SQLite no suele ser la mejor opción en sistemas de producción muy
extensos, pero para proyectos de menor envergadura su rendimiento es más que
suficiente.

\subsubsection{Gestión de tareas}

Para la gestión de tareas de forma asíncrona se ha utilizado \textbf{Celery},
una cola de tareas asíncrona enfocada a operaciones en tiempo real. Celery está
escrito en Python, pero se integra fácilmente con cualquier otro sistema que
implemente su interfaz. Las unidades de ejecución básicas, las \textit{tareas},
se ejecutan de forma concurrente en uno o varios \textit{workers}. Las tareas
pueden ejecutarse de forma asíncrona, en segundo plano, o de forma síncrona,
haciendo que el flujo de ejecución espere hasta que el resultado de la tarea
esté disponible.

Celery se comunica mediante mensajes, normalmente utilizando un \textbf{broker}
que media entre los clientes y los workers. Para iniciar una tarea, el cliente
pone un mensaje en la cola del broker. Entonces, el broker entrega el mensaje al
worker de Celery. La opción recomendada, y la que se ha utilizado en SiteUp, es
usar \textbf{RabbitMQ} como broker de mensajes.

RabbitMQ, y la mayoría de brókers de mensajes, implementan el protocolo
\ac{AMQP}, un estándar abierto para enviar mensajes entre aplicaciones y
organizaciones, siempre que éstas sean capaces de interactuar usando las
interfaces del protocolo.

\subsubsection{Dependencias de Django}

En el proyecto se han utilizado un gran número

%%% Local Variables: 
%%% mode: latex
%%% TeX-master: "../memoria"
%%% End: 


\chapter{Pruebas}
\label{chap:pruebas}
%En este capítulo se detallan las pruebas a las que se ha sometido
\textbf{oFlute}. Habitualmente, probar un videojuego o aplicación con gran
interactividad suele ser un proceso complejo. Dada la gran cantidad de
situaciones distintas que pueden tener lugar, es difícil recrear y probar todos
los escenarios posibles. Tanto es así que, en los grandes equipos de desarrollo,
hay empleados que se dedican al completo a probar los productos.

En oFlute, las pruebas se han organizado al final de cada iteración, comprobando
el correcto funcionamiento de cada nueva funcionalidad, por separado, y luego
integrando el conjunto. Además, ciertos módulos no asociados a iteraciones (como
el sistema de animaciones, sección \ref{sec:animaciones}) también fueron
probados de manera independiente.

\section{Pruebas unitarias}
Las \textbf{pruebas unitarias} se encargan de comprobar el correcto
funcionamiento de módulos y fragmentos por separado. En nuestro caso, todos los
módulos fueron probados de forma independiente, utilizando a tal efecto ramas de
desarrollo alternativas (\textit{branches}), antes de integrarlos en la rama de
desarrollo principal.

Mención especial merecen las pruebas unitarias relacionadas con el analizador
básico de notas, en la segunda iteración del proyecto. Tal y como se comentó en
la sección \ref{sec:implementacion_analizador}, al no tener experiencia previa
se tuvieron que hacer bastantes pruebas con el código relacionado con el control
del sonido para que el rendimiento fuera correcto y no apareciesen
problemas. Las pruebas unitarias en este módulo dieron lugar a decisiones de
cambios de bibliotecas en este módulo. Además, al ejecutarse en un hilo
separado, el código del analizador era propenso a generar fugas de memoria si no
se controlaban bien la vida y visibilidad de las variables utilizadas. 

\section{Pruebas de integración}
Según se comprobaban los módulos de manera individual, se iban integrando en la
rama de desarrollo principal, pasando a formar parte de la aplicación. Esta
integración también ha estado sujeta a pruebas, por un lado comprobando que los
módulos no cambiaban su comportamiento al ser integrados con el resto, y por
otro verificando que el funcionamiento general del sistema siguiera siendo
correcto.

Como se comentó en la sección \ref{sec:implementacion_estados}, uno de los
errores que más dificultad mostró, asociado a la gestión de estados, fue también
el más difícil de encontrar y reproducir. Gracias a las pruebas de integración y
al uso de herramientas de depuración, su resolución fue un éxito y sentó las
bases para un nuevo gestor de estados (sección
\ref{sec:resolucion_problema_estados}).

\section{Pruebas de jugabilidad}
Las \textbf{pruebas de jugabilidad} permitieron encontrar y rellenar lagunas en
el proyecto, haciendo que su uso fuera más sencillo para el usuario. En esta
etapa del proyecto se contó con la ayuda de usuarios reales inicialmente ajenos
al proyecto.

Entre las mejoras que surgieron a raíz de las pruebas de jugabilidad se
encuentra la implementación de la \textbf{barra de resaltado} (ver
\textit{\nameref{chap:manual_usuario}}), que aparece durante la interpretación
de las canciones para ayudar al usuario a conocer qué notas está tocando en cada
momento. Esto mejoró en gran medida las puntuaciones de los jugadores en las
partidas.

\section{Pruebas de interfaz}
Por último, una vez ajustada la jugabilidad y verificado que los módulos
funcionan bien de forma independiente y en integración, se dedicó un tiempo a
comprobar que la interfaz gráfica de usuario cumplía los requisitios mínimos que
se esperaban.

En particular, se hizo mucho hincapié en la duración de cada una de las
animaciones de cada sección. Animaciones largas podían llegar a ser tediosas,
mientras que unas animaciones cortas darían una desagradable sensación de
premura que no interesaba. Una vez ajustadas, se implementó la opción de poder
omitir las animaciones en pantalla, en cualquier sección y momento, mediante la
pulsación de la tecla \texttt{escape}. Esto permitiría que jugadores avanzados
pudieran llegar rápidamente a la sección de su interés de forma rápida.

Otra de las adiciones que surgieron tras las pruebas de interfaz fueron los
atajos de teclado. Se añadieron atajos para las opciones del menú principal, así
como el menú de selección de canciones y el de selección de lecciones. También
se habilitó la tecla \texttt{escape} para volver atrás en cualquier momento.

%%% Local Variables: 
%%% mode: latex
%%% TeX-master: "../memoria"
%%% End: 


\chapter{Conclusiones}
\label{chap:conclusiones}
%Durante el transcurso del desarrollo de oFlute, y sobre todo al término del
mismo, se han obtenido unas conclusiones y unos resultados, tanto de forma
personal como para con la comunidad, que intentaremos reflejar en este capítulo.


\section{Objetivos cumplidos}
Al término del desarrollo del proyecto, el proyecto ha completado todos los
objetivos a cumplir, presentes en el planteamiento inicial y detallados en la
sección \textit{\ref{sec:objetivos}}. En particular:

\begin{itemize}
\item Se llevó a cabo la creación del módulo de análisis de sonido, consiguiendo
  una efectividad en la detección de las notas muy alta.
\item Se hizo uso del módulo en una sección de análisis de notas y en el propio
  juego de interpretación de canciones, con resultados muy satisfactorios y
  funcionamiento correcto.
\item Además, el sistema de lecciones planteado se llevó a cabo correctamente,
  consiguiendo un motor dinámico bastante potente y fácilmente ampliable, con
  opción a añadir nuevas lecciones al sistema (apéndice
  \ref{chap:manual_lecciones}).
\item Finalmente, todos los objetivos se llevaron a cabo respetando la premisa
  de mantener una interfaz de usuario amigable y fluida, agradable a la vista y
  sencilla de usar.
\end{itemize}

\section{Conclusiones personales}

oFlute ha sido el proyecto más longevo al que me he enfrentado hasta ahora. A
pesar de tener conciencia de la envergadura del mismo desde el principio, el
tiempo para completarlo ha superado todas mis expectativas, sobre todo en lo que
a documentación se refiere. 

La causa de esto ha sido, en primer lugar, una planificación inicial algo
ineficiente y, por otro lado, mi recelo sobre los procedimientos actuales sobre
ingeniería del software. De cualquier modo, estoy bastante satisfecho con lo
obtenido, y esto me ha servido para aprender a regirme por una planificación de
forma más estricta.

\subsection{Conocimiento adquirido}

Gracias a oFlute he aprendido las técnicas básicas de la programación de audio,
sobre todo en la parte técnica más que teórica. Es esta parte teórica, sobre
todo la de análisis, la que me gustaría reforzar en un futuro, ahora que cuento
con las bases para conseguirlo. Como se comentó en el capítulo
\textit{\nameref{chap:fundamentos}}, los videojuegos relacionados con el audio
son un nicho aún poco explorado y que puede dar muchas satisfacciones, sobre
todo cuando el producto se orienta a un público joven.

Por otro lado, oFlute también me ha ayudado a aprender a usar bastantes
tecnologías auxiliares, en algunos casos meras herramientas que han facilitado
el trabajo y, en otros casos, elementos que resultaron ser de inestimable ayuda
e imprescindible uso al término del proyecto.

Una de estas tecnologías es \textbf{Boost}~\cite{boost}, un conjunto de
bibliotecas para C++ que amplían en gran medida la biblioteca estándar del
lenguaje. Un amplio número de componentes de Boost formarán parte del nuevo
estándar C++0x, por lo que me ha servido para ponerme al día en las
novedades que están por llegar.

\begin{figure}[htp!]
  \centering
  \includegraphics[width=0.4\textwidth]{8_conclusiones/imagen_logoboost}
  \caption{Logotipo de Boost}
\end{figure}

Utilizando como material de apoyo el libro \textit{``Beyond the C++ Standard
  Library''}~\cite{libroboost}, pude conocer una gran parte de las bibliotecas
de Boost, incluyendo punteros inteligentes, programación funcional, bibliotecas
para hilos, y utilizarlas en el proyecto.

Al tratarse de la principal biblioteca utilizada durante el desarrollo, oFlute
me ha provisto de un profundo conocimiento de \textbf{Gosu}~\cite{gosu},
permitiéndome implementar videojuegos con mucha más fluidez y labrándome un
pequeño \textit{framework} personal que utilizar de base en próximos
proyectos. 

\begin{figure}[htp!]
  \centering
  \includegraphics[width=0.4\textwidth]{8_conclusiones/imagen_logogosu}
  \caption{Logotipo de Gosu}
\end{figure}

Otra de las tecnologías que he aprendido ha sido \textbf{GNU Gettext}, que ha
servido para internacionalizar el proyecto. Aunque inicialmente se pensó en
utilizar una solución propia para la internacionalización el proyecto,
posteriormente se decidió usar esta tecnología, ampliamente conocida y
utilizada.

\section{Difusión y derivados}

A raíz del desarrollo de oFlute se abrieron varias vías de trabajo, en forma de
actividades de difusión de conocimiento, proyectos informáticos y participación
en comunidad.

\subsection{Conocimiento generado}

Además de servir como recurso para el correcto transcurso del proyecto, los
conocimientos adquiridos me permitieron, en numerosos casos, generar
documentación adicional e impartir talleres relacionados las tecnologías
previamente mencionadas.

En cuanto a Boost, se llevó a cabo un taller durante los \textit{Cursos de
  Verano de la OSLUCA}~\cite{cursosverano}, en el que se explicaron las partes
más importantes de esta colección de bibliotecas, con numerosos ejemplos de
muchas de ellas. Toda la documentación es libre~\cite{materialesCursoBoost}.

Por otro lado, impartí un taller~\cite{tallergosu} sobre la biblioteca Gosu en
colaboración con la ADVUCA~\cite{advuca}, cuya afluencia superó las 50
personas. Los materiales pueden descargarse
libremente~\cite{tallergosumateriales}.

Por último, tras la aplicación de GNU Gettext en oFlute, escribí una guía
concisa sobre traducción de proyectos con esta herramienta, que se puede
encontrar en el apéndice~\textit{\nameref{chap:tutorial_gettext}}.


\subsection{Freegemas}

Uno de estos proyectos \textit{``hijos''} de oFlute ha sido
\textbf{Freegemas}~\cite{freegemas}, un clon libre del popular juego tipo puzzle
\textit{Bejeweled}. Freegemas es multiplataforma, funciona tanto en GNU/Linux
como en Windows, y además forma parte oficial de Guadalinex~\cite{guadalinex},
por lo que es posible encontrarlo en los repositorios oficiales.

\begin{figure}[h!]
  \centering
  \includegraphics[width=0.55\textwidth]{8_conclusiones/imagen_freegemas}
  \caption{Logotipo de Freegemas}
\end{figure}

Además, Freegemas sirvió como base para una serie de tres artículos que
publiqué, junto al director del presente PFC, en la revista Linux
Magazine~\cite{linuxmagazine} sobre desarrollo de videojuegos en C++. Es posible
encontrar estos artículos en el archivo de la
revista~\cite{refarticulo1}~\cite{refarticulo2}~\cite{refarticulo3} bajo una
licencia Creative Commons.

\subsection{IV Concurso Universitario de Software Libre}

Por otro lado, gracias a oFlute tuve la oportunidad de participar en el IV
Concurso Universitario de Software Libre~\cite{cusl}. En el transcurso del
concurso formé parte de una comunidad muy unida, en la que reinó el apoyo y la
ayuda entre los concursantes. La final del concurso se celebró en la Escuela
Superior de Ingeniería de Cádiz, en la que oFlute obtuvo una mención
especial~\cite{cusl2}.

\begin{figure}[h!]
  \centering
  \includegraphics[width=0.65\textwidth]{8_conclusiones/imagen_logocusl}
  \caption{IV Concurso Universitario de Software Libre}
\end{figure}

Previa a la final nacional tuvo lugar la fase local del concurso, en el que el
proyecto también recibió un accésit al mejor proyecto de
innovación~\cite{cusllocal}.

%%% Local Variables: 
%%% mode: latex
%%% TeX-master: "../memoria"
%%% End: 


\appendix

\chapter{Herramientas utilizadas}
\label{chap:herramientas}
%En este primer apéndice se hará un repaso por las herramientas que han hecho
posible el desarrollo de oFlute, indicando qué tareas se han llevado a cabo con
cada una de ellas.

Dado que en la actualidad, la tendencia es la de mover la potencia de cálculo a
la \textit{nube}, también se mencionarán las herramientas \textit{online} que
han aportado al proyecto.

\section{Lenguajes y bibliotecas de programación}

\subsection{C++}
El lenguaje elegido para elaborar el proyecto fue \textbf{C++}. C++ es un
lenguaje de programación que extendió el lenguaje original C hacia el paradigma
de la orientación a objetos, mediante el uso de clases, y el paradigma de la
programación genérica, mediante plantillas, además de incluir muchas otras
novedades y facilidades.

Se decidió utilizar C++ por la familiaridad que teníamos con el mismo, así como
por la velocidad y optimización que se obtienen en comparación con otros
lenguajes, como Java.

\subsection{Gosu}
\textbf{Gosu}~\cite{gosu} es una biblioteca libre para el desarrollo de
videojuegos 2D para Ruby y C++. Entre sus características más atractivas se
encuentran la aceleración gráfica por hardware, la sencillez de su API,
completamente orientada a objetos, y su compatibilidad con sistemas GNU/Linux,
Windows y Mac OS.

Esta biblioteca, que está cerca de lanzar su versión 0.8, cuenta con una base de
usuarios cada vez mayor, en parte gracias a su relativa popularidad entre los
desarrolladores indie y en los eventos de programación rápida de videojuegos.

\subsection{PulseAudio}
\textbf{PulseAudio}~\cite{pulseaudio} (mencionado previamente en la sección
\ref{sec:pulseaudio1}) es un servidor de sonido multiplataforma, muy popular en
distribuciones GNU/Linux actuales. Se decidió su uso por lo intuitiva y sencilla
de utilizar que resultaba su API, así como su estabilidad, en comparaciones con
otras opciones previamente probadas.

\subsection{GNU Gettext}

\textbf{GNU Gettext}~\cite{refgettext} es un conjunto de herramientas de internacionalización de
proyectos muy sencilla y popular. Se utilizó porque es la solución más eficaz y
aceptada a la hora de traducir un proyecto. Se puede encontrar un manual
introductorio en el apéndice \textit{\nameref{chap:tutorial_gettext}}.

\subsection{KissFFT}
\textbf{KissFFT}~\cite{kissfft} es una biblioteca que implementa el algoritmo
FFT para el cálculo de transformadas rápidas de Fourier. Es una biblioteca
pequeña, razonablemente eficiente y portable, con capacidad para realizar
operaciones en distintos formatos.

Dado que el módulo de análisis de sonidos se basa en la transformada de Fourier,
era necesario encontrar una biblioteca que pudiera realizar el cálculo de forma
rápida y precisa. KissFFT resultó ser una candidata ideal para esta tarea.

\subsection{PugiXML}
\textbf{PugiXML}~\cite{pugixml} es una ligera biblioteca para el parseo y
procesamiento de ficheros XML. Tiene soporte completo Unicode, un parseador muy
veloz y capacidad para usar consultas XPath. Además, su implementación apenas
ocupa cuatro ficheros de trivial compilación, por lo que resulta muy sencillo
incluirla en cualquier proyecto.

Esta biblioteca se utilizó para el procesamiento de los ficheros de canciones y
lecciones, que utilizan XML para representar los distintos elementos.

\subsection{Boost}
\textbf{Boost}~\cite{boost} es un conjunto de bibliotecas para C++ para un gran
cantidad de situaciones distintas. Entre sus componentes se incluyen bibliotecas
para el procesamiento de textos, herramientas de gestión de memoria, expresiones
regulares, parseo de opciones de línea de comandos, adaptadores para
programación funcional y un largo etcétera.

Entre sus desarrolladores se encuentran muchos de los mejores programadores de
C++, y tal es su calidad y popularidad, que 10 de las bibliotecas que componen
Boost formarán parte del nuevo estándar de C++.

En oFlute se han utilizado bastantes partes de Boost: punteros inteligentes,
soporte para expresiones regulares, acceso al sistema de ficheros, y operaciones
con cadenas. Esto ha simplificado mucho el código, consiguiendo un estilo más
limpio y elegante, sin sacrificar funcionalidad alguna.

\section{Gestión de código}

\subsection{Subversion}
\textbf{Subversion}~\cite{refsubversion}, popularmente conocido como
\textbf{SVN}, es un sistema de control de versiones muy popular y
utilizado. Hasta hace unos años era el indiscutible ganador entre los sistemas
de control de versiones de código, y en la actualidad sigue siendo la opción
principal en muchas de las forjas de código más importantes de internet.

Se basa en un repositorio central, al que se envían los cambios de
los ficheros versionados. Además, a diferencia de otros sistemas como CVS, lleva
un control de revisiones a nivel global, no a nivel de fichero.

Subversion ha servido, en oFlute, sobre todo como backup online, además de
historial de cambios. En más de una ocasión ha sido necesario revertir los
cambios, y Subversion ha facilitado mucho el proceso, que en otras
circunstancias podría no haber sido posible.

\subsection{Forjas: RedIris y Google Code}
Para llevar un control del código del proyecto y servir como punto de encuentro
con otros usuarios y desarrolladores se abrió una forja para oFlute,
inicialmente dentro de la \textbf{Forja de Conocimiento Libre de la Comunidad
  RedIRIS}\footnote{\url{http://forja.rediris.es}}, gestionada por la Junta de
Andalucía. Esta forja ofrecía, además de un sistema de control de versiones del
código, una serie de herramientas para la gestión de bugs, emisión de noticias,
listas de correos, etcétera.

Desafortunadamente la forja de RedIris carecía de un mantenimiento adecuado, lo
que llevó al traspaso del proyecto a \textbf{Google Code}~\cite{ofluteforja}, un
repositorio de código ofrecido por Google mucho más actualizado y con un
mantenimiento adecuado. Una de los atractivos de Google Code es su sistema de
páginas wiki, que permite gestionar documentación online utilizando una sintaxis
sencilla y potente.

\subsection{Doxygen}
\textbf{Doxygen}~\cite{doxygen} es una herramienta de generación automática de
documentación de código fuente. Basándose en unos comentarios con una sintaxis
especial, Doxygen genera documentación en un gran número de formatos: HTML, PDF,
RTF, \LaTeX...

Esta herramienta se utilizó para generar toda la documentación del código del
proyecto en formato \ac{HTML}.


\section{Herramientas de diseño y diagramas}

\subsection{Adobe Photoshop}
\textbf{Adobe Photoshop}~\cite{photoshop} es la herramienta de edición gráfica de mapa de bits
más popular y utilizada en el mundo. Su longeva historia y profuso desarrollo ha
hecho que se convierta en la herramienta estándar a la hora de trabajar con
imágenes no vectoriales.

Aunque no se trata de una herramienta libre, se utilizó esta aplicación para la
creación de todo el apartado visual de oFlute. Todas las secciones fueron
diseñadas con esta herramienta, obteniendo un apartado gráfico atractivo y
minimalista.

\subsection{Inkscape}
\textbf{Inkscape}~\cite{inkscape} es una herramienta libre de edición de gráficos vectoriales. A
diferencia de Adobe Photoshop, Inkscape trabaja con formas definidas
matemáticamente, que es posible escalar y distorsionar sin causar pérdida de
datos. Esto hace que sea la herramienta perfecta para trabajar con logotipos y
diagramas.

Esta aplicación ayudó al diseño del logotipo oficial de oFlute, así como al
ajuste de ciertos diagramas pertenecientes a la documentación.

\subsection{Dia}
\textbf{Dia}~\cite{dia} es un editor de diagramas libre, parte de la suite ofimática de
GNOME. Tiene un diseño modular que le permite abordar una gran variedad de
diagramas distintos, desde diagramas de flujo hasta modelos entidad-relación,
UML, etcétera.

Este software se empleó para la generación de la mayor parte de los diagramas
presentes en esta memoria.

\subsection{BoUML}

\textbf{BOUML}~\cite{bouml} es una herramienta libre para diseñar diagramas siguiendo la
notación UML. Cuenta con utilidades de generación automática de código a partir
de diagramas en lenguajes C++, Java, PHP, Python e IDL. Es multiplataforma y
permite la generación de varios tipos de diagramas: diagramas de clases, de
interacción, de secuencia, de casos de uso, etc

\subsection{ImageMagick}

\textbf{ImageMagick}~\cite{imagemagick} es un conjunto de utilidades de línea de comandos para
mostrar, editar y convertir ficheros de imagen. Resultan especialmente útil a la
hora de procesar un gran número de imágenes de forma automática, como por
ejemplo cuando es necesario convertir varias imágenes de un formato a otro, o
cambiar el tamaño de numerosos ficheros.

\section{Edición de textos}

\subsection{GNU Emacs}
\textbf{GNU Emacs}~\cite{refemacs} es un potente editor multiplataforma. Su propio manual lo
describe como \textit{``un editor extensible, personalizable, auto-documentado y
  de tiempo real''}. Emacs es tan potente que a menudo los usuarios
experimentados no tienen que salir de su interfaz para realizar cualquier
operación, ya que desde el propio editor se puede, desde navegar por la red,
hasta revisar el correo electrónico.

Emacs ha sido el editor utilizado para escribir todo el código fuente del
proyecto, así como toda la documentación en \LaTeX. Sus diversos \textit{modos}
se ajustan a cada tipo de documento, proporcionando funciones que faciliten y
aceleren la edición de los textos.

\subsection{\LaTeX}
\textbf{\LaTeX}~\cite{latex} es un sistema de composición de textos, orientado especialmente
a la creación de libros, documentos científicos y técnicos que contengan
fórmulas matemáticas. Es muy utilizado para la composición de artículos
académicos, tesis y libros técnicos, dado que la calidad tipográfica de los
documentos realizados con LaTeX es comparable a la de una editorial científica
de primera línea.

La presente memoria ha sido escrita y maquetada con \LaTeX.

%%% Local Variables: 
%%% mode: latex
%%% TeX-master: "../memoria"
%%% End: 


\chapter{Manual de instalación}
\label{chap:manual_instalacion}
%A continuación se presentan las instrucciones de instalación de \textbf{oFlute}
en sistemas GNU/Linux basados en paquetería Debian~\cite{refdebian}, por
ser los que más auge están teniendo en la actualidad. Tomaremos como referencia
la distribución Ubuntu 10.10 Maverick~\cite{refubuntu}.

\section{Instalación de dependencias}

El primer paso a la hora de instalar oFlute es hacernos con las dependencias del
proyecto. Para ello, utilizaremos el gestor \texttt{apt-get}, pero para otros
sistemas el procedimiento será similar.

Primero, instalaremos las dependencias de la biblioteca Gosu:

\begin{minted}{bash}
  sudo apt-get install g++ libgl1-mesa-dev libpango1.0-dev \
  libopenal-dev libsndfile-dev libxdamage-dev libsdl-ttf2.0-dev \
  libboost-dev libfreeimage3 libfreeimage-dev
\end{minted}

Seguidamente, instalaremos el resto de dependencias:

\begin{minted}{bash}
  sudo apt-get install libboost-regex-dev libboost-filesystem-dev \ 
  libboost-thread-dev libpulse-dev
\end{minted}

\section{Descarga del código fuente}

El siguiente paso será hacernos con una copia local del repositorio, que
contiene el código fuente. Al tratarse de un proyecto libre, el repositorio es
de acceso público~\cite{ofluteforja}.

Haciendo uso del comando \texttt{svn}, haremos una copia local del repositorio
Subversion~\cite{refsubversion} en nuestro sistema con la siguiente orden:

\begin{minted}{bash}
  svn checkout http://oflute.googlecode.com/svn/trunk oflute
\end{minted}

Con ello, se creará una carpeta llamada \texttt{oflute} en la que se encontrarán
los ficheros de la rama principal (\texttt{trunk}) del proyecto.

\section{Compilación}

Como se ha comentado previamente, será necesario compilar la biblioteca Gosu
previamente. Para ello, se ha facilitado un objetivo en el script de compilación
\texttt{Makefile}~\cite{refmakefile} que automatiza el proceso.

Así, el primer paso será, estando en la carpeta recién creada, utilizar el
siguiente comando:

\begin{minted}{bash}
  make libgosu
\end{minted}

Una vez concluida la compilación de Gosu, pasaremos a compilar el propio
proyecto mediante la siguiente orden:

\begin{minted}{bash}
  make
\end{minted}

\section{Configuración del micrófono}
Antes de poder lanzar el juego, será necesario configurar las opciones de audio
del sistema operativo. Por defecto, la entrada de audio suele estar mutada, por
lo que será imposible que la aplicación tenga acceso al sonido proveniente del
micrófono.

Siguiendo con la referencia de Ubuntu, iremos al menú \textit{Sistema}, luego a
\textit{Preferencias} y, finalmente, elegiremos \textit{Sonido}. Una vez ahí,
pasaremos a la pestaña \textit{Entrada} y, en caso de que estuviera marcada,
desmarcaremos la opción \textit{Silenciar}, tal y como muestra la imagen.

\begin{figure}[h!]
  \centering
  \includegraphics[width=0.75\textwidth]{apendice_manual_instalacion/imagen_captura_1}
  \caption{Panel de configuración del micrófono}
\end{figure}

\vspace{-1cm}

\section{Ejecución}

Al término de las anteriores etapas, se habrá generado un ejecutable de nombre
\texttt{oflute}, por lo que será posible lanzar la aplicación a través del
siguiente comando:

\begin{minted}{bash}
  ./oflute
\end{minted}

%%% Local Variables: 
%%% mode: latex
%%% TeX-master: "../memoria"
%%% End: 


\chapter{Manual de usuario}
\label{chap:manual_usuario}
%En este capítulo se explicará el funcionamiento de la aplicación desde el punto
de vista del usuario final, detallando las opciones de cada una de las secciones.

Para poder acceder a la aplicación, tal y como se explicó en el apéndice 
\textit{\nameref{chap:manual_instalacion}}, será necesario utilizar el comando:

\begin{minted}{bash}
  ./oflute
\end{minted}

\section{Ejecución inicial}
Al lanzar la aplicación, se abrirá una ventana y aparecerán, consecutivamente,
la pantalla de créditos del autor y la pantalla de presentación del
juego. Pasados unos momentos éstas desaparecerán, aunque tiene la opción de
omitir cada una de las pantallas pulsando la tecla \texttt{escape}.

\begin{figure}[h!]
  \centering
  \includegraphics[width=0.8\textwidth]{apendice_manual_usuario/imagen_estadoAutor}
  \caption{Pantalla de créditos del autor}
  \vspace{-1cm}
\end{figure}

\begin{figure}[h!]
  \centering
  \includegraphics[width=0.8\textwidth]{apendice_manual_usuario/imagen_estadoIntro}
  \caption{Pantalla de presentación del juego}
  \vspace{-0.3cm}
\end{figure}

Una vez desaparezcan ambas pantallas, se iniciará una animación para mostrar el
menú principal. Puede cancelar la animación y mostrar rápidamente el menú
principal pulsando la tecla \texttt{escape}.

\begin{figure}[h!]
  \vspace{-0.1cm}
  \centering
  \includegraphics[width=0.8\textwidth]{apendice_manual_usuario/imagen_menuPrincipal}
  \caption{Pantalla del menú principal}
  \vspace{-1cm}
\end{figure}

Cuando las animaciones concluyan, el menú principal ya será completamente
funcional. Podremos acceder a las diferentes opciones, ya sea haciendo click con
el ratón en los botones de pantalla, o pulsando alguna de las teclas de acceso
directo. Las secciones y teclas asignadas son las siguientes:

\begin{itemize}
\item \textbf{Analizador de notas} (\textit{tecla 1}) Abre el analizador de
  notas.
\item \textbf{Canciones} (\textit{tecla 2}) Pasa al menú de selección de
  canciones.
\item \textbf{Lecciones} (\textit{tecla 3}) Abre el menú de elección de
  lecciones.
\item \textbf{Calibrar micrófono} (\textit{tecla 4}) Lanza el sistema de
  calibración del micrófono.
\item \textbf{Salir} (\textit{tecla 5}) Cierra la aplicación
\end{itemize}

Lo recomendado, en la primera ejecución de \textbf{oFlute}, es lanzar la sección
\textit{Calibrar micrófono}, ya sea pulsando el botón con el ratón o mediante la
tecla 4. Con esto nos cercioramos de que el programa es capaz de distinguir el
sonido del micrófono del ruido de fondo. Una vez hecho, ya estará todo listo
para acceder al resto de secciones.

\section{Sección -- calibrar micrófono}
Pulsando el botón correspondiente, o mediante la tecla 4, pasaremos a la sección
de calibración del micrófono. En primer lugar, aparecerá el siguiente mensaje:

\begin{figure}[h!]
  \vspace{-0.1cm}
  \centering
  \includegraphics[width=0.8\textwidth]{apendice_manual_usuario/imagen_seccionCalibrar1}
  \caption{Pantalla inicial de calibración del micrófono}
  \vspace{-1cm}
\end{figure}

Como se indica en el mensaje, para comenzar la calilbración el usuario deberá
pulsar la tecla \texttt{espacio} y guardar silencio en el micrófono durante dos
segundos. En ese momento, aparecerá el siguiente mensaje:

\begin{figure}[h!]
  \vspace{-0.1cm}
  \centering
  \includegraphics[width=0.8\textwidth]{apendice_manual_usuario/imagen_seccionCalibrar_mensaje1}
  \caption{Mensaje al inicio del calibrado}
\end{figure}

Una vez pasen los dos segundos y concluya la calibración, se mostrará el
siguiente mensaje:

\begin{figure}[h!]
  \vspace{-0.1cm}
  \centering
  \includegraphics[width=0.8\textwidth]{apendice_manual_usuario/imagen_seccionCalibrar_mensaje2}
  \caption{Mensaje al final del calibrado}
\end{figure}

Cuando concluya el proceso, la aplicación habrá guardado el valor umbral de
ruido del micrófono, lo que facilitará el análisis del sonido de la flauta.

Hay casos en los que la calibración será defectuosa, y un mensaje informará de
ello. En tal caso, no habrá más que comprobar la configuración del micrófono y
repetir el proceso de calibrado.

\section{Sección -- analizador de notas}

Al pulsar el botón correspondiente a la sección del analizador de notas,
mediante una animación aparecerán los elementos de esta parte de la aplicación
(figura \ref{fig:pantalla_analizador_notas} en página
\pageref{fig:pantalla_analizador_notas}).

\begin{figure}[h!]
  \vspace{-0.1cm}
  \centering
  \includegraphics[width=0.75\textwidth]{apendice_manual_usuario/imagen_seccionAnalizador}
  \caption{Pantalla del analizador de notas}
  \label{fig:pantalla_analizador_notas}
\end{figure}

Una vez terminen de mostrarse todos los elementos, el usuario podrá empezar a
utilizar la funcionalidad de la sección. En concreto, el analizador de notas nos
permite ver reflejada en el pentagrama la nota que toquemos con la flauta.

Su uso, pues, es muy sencillo. Simplemente tendremos que tocar nuestra flauta de
forma que el micrófono sea capaz de captar su sonido. Si lo hacemos
correctamente, el analizador mostrará en cada momento la nota que está tocando
sobre el pentagrama.

Una vez hayamos acabado de utilizar esta sección de la aplicación, podremos
pulsar en el botón \textit{Volver}, o en la tecla \texttt{escape} del teclado,
lo que nos llevará de vuelta al menú principal.

\section{Sección -- lecciones}
\vspace{-0.5cm}
Si seleccionamos la opción \textit{Lecciones} en el menú principal, o
alternativamente presionamos la tecla 3 del teclado, la aplicación mostrará una
animación para ocultar el menú principal y nos mostrará el menú de selección de
lecciones. 

\begin{figure}[h!]
  \vspace{-0.2cm}
  \centering
  \includegraphics[width=0.75\textwidth]{apendice_manual_usuario/imagen_seccionLecciones1}
  \caption{Pantalla del menú de selección de lecciones}
  \vspace{-1cm}
\end{figure}

\pagebreak

En cualquier momento podemos hacer click en el botón \textit{Volver al menú}, o
pulsar la tecla \texttt{escape} del teclado, para ir a la sección anterior.

En el menú de selección de lecciones podremos encontrar los siguientes
elementos:

\begin{itemize}
\item \textbf{Título} de la lección
\item \textbf{Descripción} de la lección
\item Botón \textbf{Comenzar lección}.
\item Botón \textbf{Anterior lección}.
\item Botón \textbf{Siguiente lección}.
\item Botón \textbf{Volver al menú}.
\end{itemize}

Mediante los botones de \textit{anterior} y \textit{siguiente lección} podremos
navegar entre las diferentes lecciones cargadas en el sistema. Haciendo uso de
los paneles de \textit{título} y \textit{descripción} nos informaremos sobre la
lección elegida. Una vez que tengamos claro la lección que queremos tomar,
pulsaremos en el botón \textit{comenzar lección.}

Una vez que decidamos comenzar la lección, el sistema ocultará los elementos del
menú de selección de lecciones mediante animaciones y, posteriormente, cargará y
mostrará los elementos de la lección elegida.

\begin{figure}[h!]
  \centering
  \includegraphics[width=0.8\textwidth]{apendice_manual_usuario/imagen_seccionLecciones2}
  \caption{Pantalla de lección de ejemplo}
\end{figure}

Las lecciones se muestran en forma de conjunto de elementos multimedia --
imágenes y texto -- de fácil comprensión, que el usuario podrá leer y estudiar
de forma independiente. En futuras versiones de la aplicación se podrán utilizar
lecciones de varias etapas, así como integrar sonidos y otro tipo de multimedia.

\section{Sección -- canciones}

El usuario podrá acceder a esta sección desde el menú principal, pulsando en la
botón \textit{Canciones}, o con la tecla 2 del teclado.

Al acceder, desaparecerá el menú principal y se mostrará el menú de selección de
canciones.

\begin{figure}[h!]
  \centering
  \includegraphics[width=0.8\textwidth]{apendice_manual_usuario/imagen_seccionCanciones1}
  \caption{Pantalla del menú de selección de canciones}
\end{figure}

Esta pantalla consta de los siguientes botones:
\begin{itemize}
\item \textbf{Anterior canción}, representado con una flecha hacia arriba.
\item \textbf{Siguiente canción}, representado con una flecha hacia abajo.
\item \textbf{Comenzar canción}, simbolizado con la palabra \textit{OK}.
\item \textbf{Volver}.
\end{itemize}

Mediante los botones \textit{anterior} y \textit{siguiente canción}, el usuario
podrá elegir el tema a interpretar con la flauta. Una vez esté resaltada la
canción correcta, el usuario deberá pulsar el botón \textit{Comenzar canción --
  OK} para que dé comienzo la interpretación.

También es posible pulsar el botón \textit{volver} para ir de nuevo al menú
principal.

\subsection{Interpretación de la canción}

Al lanzar la canción, aparecerá la pantalla de interpretación de canción, que
contiene numerosos elementos importantes para el jugador.

\begin{figure}[h!]
  \centering
  \includegraphics[width=0.8\textwidth]{apendice_manual_usuario/imagen_seccionCanciones2}
  \caption{Pantalla de interpretación de canción}
\end{figure}

\begin{enumerate}
\item \textbf{Marcador} de la puntuación del usuario.
\item \textbf{Pentagrama}. Sobre él aparecerán las notas.
\item \textbf{Notas}. Éstas irán apareciendo por el lado derecho del pentagrama,
  y moviéndose hacia la izquierda de forma ordenada y rítmica.
\item \textbf{Resaltado de nota}\label{barra_resaltado}. Resalta la posición en
  el pentagrama de la nota que está tocando el usuario con la flauta en cada
  momento. Esto ayuda visualmente a saber si estamos interpretando la nota
  correcta o no.
\item \textbf{Barra de interpretación}. Esta barra indica cuándo debe el usuario
  empezar a tocar la nota. 
\item \textbf{Barra de progreso}. Indica el progreso de la interpretación de la
  canción. Cuando la barra esté completa, la canción concluirá.
\end{enumerate}

La forma de juego es sencilla. Como se ha indicado, las notas aparecerán sobre
el pentagrama por el lado derecho, e irán desplazándose hacia la izquierda. Una
vez que una nota llegue a la \textit{barra de interpretación}, el jugador deberá
tocar con la flauta la nota correcta, de forma constante hasta que llegue el
turno de la siguiente nota.

Para ayudar al jugador, la barra de \textit{resaltado de nota} indicará en cada
momento qué nota se está tocandoo. Así, si en un instante es necesario tocar un
\textit{Sol} pero la barra de resaltado está por encima, el usuario sabrá que
debe cerrar más orificios de la flauta.

Las posibles \textbf{notas} que pueden aparecer son \textit{Do}, \textit{Re},
\textit{Mi}, \textit{Fa}, \textit{Sol}, \textit{La}, \textit{Si}, \textit{Do
  grave} y \textit{Re grave}. Las posibles figuras que podrán aparecer son las
siguientes:

\begin{description}
\item[Redonda] -- Dura cuatro tiempos.
\vspace{-0.1cm}
\begin{figure}[h!]
  \centering
  \includegraphics[width=0.05\textwidth]{apendice_manual_usuario/imagen_figRedonda}
  \caption{Figura musical -- Redonda}
\end{figure}

\vspace{-0.35cm}

\item[Blanca] -- Dura dos tiempos.
\vspace{-0.1cm}
\begin{figure}[h!]
  \centering
  \subfloat[Posición normal]{\parbox{0.5\textwidth}{\centering\includegraphics[width=1cm]{apendice_manual_usuario/imagen_figBlanca}}}
  \subfloat[Posición invertida]{\parbox{0.5\textwidth}{\centering\includegraphics[width=1cm]{apendice_manual_usuario/imagen_figBlancaInv}}}
  \caption{Figura musical -- Blanca}
\end{figure}

\vspace{-0.35cm}

\item[Negra] -- Dura un tiempo.
\vspace{-0.1cm}
\begin{figure}[h!]
  \centering
  \subfloat[Posición normal]{\parbox{0.5\textwidth}{\centering\includegraphics[width=1cm]{apendice_manual_usuario/imagen_figNegra}}}
  \subfloat[Posición invertida]{\parbox{0.5\textwidth}{\centering\includegraphics[width=1cm]{apendice_manual_usuario/imagen_figNegraInv}}}
  \caption{Figura musical -- Negra}
\end{figure}

\vspace{-0.35cm}

\item[Corchea] -- Dura la mitad que una negra.
\vspace{-0.1cm}
\begin{figure}[h!]
  \centering
  \subfloat[Posición normal]{\parbox{0.5\textwidth}{\centering\includegraphics[width=1cm]{apendice_manual_usuario/imagen_figCorchea}}}
  \subfloat[Posición invertida]{\parbox{0.5\textwidth}{\centering\includegraphics[width=1cm]{apendice_manual_usuario/imagen_figCorcheaInv}}}
  \caption{Figura musical -- Corchea}
\end{figure}
\end{description}

\vspace{-0.35cm}

De manera excepcional, es posible alargar el tiempo de una figura mediante el
uso del \textbf{puntillo}, que hará que su duración aumente la mitad de su
tiempo original. El puntillo se representa mediante un pequeño círculo a la
derecha de la base de la nota.

Por otro lado, las siguientes son las figuras que representan los
\textbf{silencios}. Funcionan igual que las notas normales, solo que el jugador,
en lugar de tocar una nota, deberá mantenerse en silencio desde que el silencio
toque la barra de interpretación hasta que llegue la siguiente figura.

\begin{description}
\item[Silencio de redonda] -- Dura cuatro tiempos, y se coloca unido, por
  debajo, a la cuarta línea del pentagrama.

\begin{figure}[h!]
  \centering
  \includegraphics[width=0.05\textwidth]{apendice_manual_usuario/imagen_silBlanca}
  \caption{Silencio de redonda}
\end{figure}


\item[Silencio de blanca] -- Dura dos tiempos, y se coloca posado sobre la
  tercera línea del pentagrama.

\begin{figure}[h!]
  \centering
  \includegraphics[width=0.05\textwidth]{apendice_manual_usuario/imagen_silBlanca}
  \caption{Silencio de blanca}
\end{figure}


\item[Silencio de negra] -- Dura lo mismo que una negra.
\begin{figure}[h!]
  \centering
  \includegraphics[width=0.05\textwidth]{apendice_manual_usuario/imagen_silNegra}
  \caption{Silencio de negra}
\end{figure}

\item[Silencio de corchea] -- Dura la mitad que una negra.
\begin{figure}[h!]
  \centering
  \includegraphics[width=0.05\textwidth]{apendice_manual_usuario/imagen_silCorchea}
  \caption{Silencio de corchea}
\end{figure}
\end{description}

El final de la interpretación vendrá indicado por la doble barra final. Una vez
que ésta llegue a la zona de interpretación, se dará por concluida la canción.

\begin{figure}[h!]
  \centering
  \includegraphics[width=0.02\textwidth]{apendice_manual_usuario/imagen_figFinal}
  \caption{Doble barra de final de pentagrama}
\end{figure}

\subsection{Resultados de la interpretación}
Tras la interpretación de la pieza, aparecerá la zona de puntuaciones.

\begin{figure}[h!]
  \centering
  \includegraphics[width=0.8\textwidth]{apendice_manual_usuario/imagen_seccionCanciones3}
  \caption{Pantalla de resultados tras la interpretación}
\end{figure}

En esta pantalla, además del \textbf{título} y \textbf{subtítulo} de la canción,
se mostrará el \textbf{porcentaje de aciertos} que ha tenido el jugador en su
interpretación. Además, aparecerá una frase de apoyo que variará según el éxito
de la partida.

Una vez llegados a este punto, el jugador puede volver a la pantalla anterior
pulsando la tecla \texttt{escape}.

%%% Local Variables: 
%%% mode: latex
%%% TeX-master: "../memoria"
%%% End: 


\chapter{GNU Free Documentation License}
\label{sec:fdl}
\input{0_misc/fdl.tex}

\bibliographystyle{paquetes/hispa-annote}
\bibliography{0_misc/bibliografia}

\end{document}

%%% Local Variables:
%%% mode: latex
%%% TeX-master: t
%%% End:
